\documentclass[12pt,openright,twoside,a4paper,english,brazil]{memoir}

% --- Codificação, Fontes e Idioma ---
\usepackage[utf8]{inputenc}
\usepackage[T1]{fontenc}
\usepackage{lmodern}
\usepackage{amsmath,amssymb}
\usepackage{babel}
\usepackage{microtype}
\usepackage[htt]{hyphenat}

% --- Ajustes de Página e Margens (Memoir) ---
\settrimmedsize{297mm}{210mm}{*}
\settypeblocksize{230mm}{155mm}{*}
\setlrmarginsandblock{30mm}{25mm}{*}
\setulmarginsandblock{30mm}{25mm}{*}
\setheadfoot{30pt}{25pt}
\checkandfixthelayout

% --- Cores e Gráficos ---
\usepackage{xcolor}
\definecolor{primary}{RGB}{16, 54, 92}        % Azul marinho profundo institucional
\definecolor{secondary}{RGB}{35, 110, 160}    % Azul médio técnico
\definecolor{accent}{RGB}{185, 75, 45}        % Terracota pedagógico
\definecolor{neutraldark}{RGB}{40, 44, 52}    % Cinza carvão para texto e códigos
\definecolor{boxbg}{RGB}{246, 249, 252}       % Fundo suave de destaque
\definecolor{boxborder}{RGB}{190, 210, 230}   % Borda de destaque azulada
\definecolor{daborange}{RGB}{253, 246, 227}   % Fundo de atenção
\definecolor{borderorange}{RGB}{230, 160, 90} % Borda de atenção
\definecolor{codebg}{RGB}{247, 248, 250}      % Fundo para código

% --- Tabelas e Listas ---
\usepackage{booktabs}
\usepackage{tabularx}
\usepackage{longtable}
\usepackage{array}
\usepackage{enumitem}
\setlist[itemize]{leftmargin=1.5em, itemsep=2pt, parsep=0pt}
\setlist[enumerate]{leftmargin=1.5em, itemsep=2pt, parsep=0pt}

% --- Caixas de Texto (tcolorbox) ---
\usepackage{tcolorbox}
\tcbuselibrary{skins,breakable}

% Caixa para Fundamentos / Conceitos Chave
\newtcolorbox{conceitochave}[1][]{%
  enhanced,
  breakable,
  colback=boxbg,
  colframe=primary,
  fonttitle=\bfseries\color{white},
  coltitle=white,
  coltext=neutraldark,
  boxrule=1pt,
  arc=3mm,
  left=8pt,
  right=8pt,
  top=6pt,
  bottom=6pt,
  title={Conceito Pedag{\'o}gico Chave},
  #1
}

% Caixa para Dicas Práticas do OpenCode
\newtcolorbox{dicaopencode}[1][]{%
  enhanced,
  breakable,
  colback=boxbg,
  colframe=secondary,
  fonttitle=\bfseries\color{white},
  coltitle=white,
  coltext=neutraldark,
  boxrule=0.8pt,
  arc=3mm,
  left=8pt,
  right=8pt,
  top=6pt,
  bottom=6pt,
  title={Dica Pr{\'a}tica OpenCode},
  #1
}

% Caixa de Objetivos de Aprendizagem do Capítulo
\newtcolorbox{objetivos}[1][]{%
  enhanced,
  breakable,
  colback=boxbg,
  colframe=accent,
  fonttitle=\bfseries\color{white},
  coltitle=white,
  coltext=neutraldark,
  boxrule=0.8pt,
  arc=3mm,
  left=8pt,
  right=8pt,
  top=6pt,
  bottom=6pt,
  title={Ao final deste cap{\'i}tulo, voc{\^e} ser{\'a} capaz de:},
  #1
}

% Caixa de Atenção / Cuidado Pedagógico
\newtcolorbox{atencaopedagogica}[1][]{%
  enhanced,
  breakable,
  colback=daborange,
  colframe=borderorange,
  fonttitle=\bfseries\color{white},
  coltitle=white,
  coltext=neutraldark,
  boxrule=1pt,
  arc=3mm,
  left=8pt,
  right=8pt,
  top=6pt,
  bottom=6pt,
  title={Aten{\c c}{\~a}o \& Integridade {\'E}tica},
  #1
}

% --- Formatação de Código-fonte (listings) ---
\usepackage{listings}
\lstset{
  backgroundcolor=\color{codebg},
  basicstyle=\ttfamily\footnotesize\color{neutraldark},
  keywordstyle=\bfseries\color{primary},
  stringstyle=\color{accent},
  commentstyle=\itshape\color{secondary},
  breaklines=true,
  frame=single,
  rulecolor=\color{boxborder},
  numbers=none,
  tabsize=2,
  showstringspaces=false,
  captionpos=b
}

% --- Estilo dos Capítulos no Memoir ---
\makechapterstyle{manualtheme}{%
  \setlength{\beforechapskip}{20pt}
  \setlength{\afterchapskip}{25pt}
  \renewcommand*{\chapnamefont}{\normalfont\Large\scshape\color{secondary}}
  \renewcommand*{\chapnumfont}{\normalfont\HUGE\bfseries\color{primary}}
  \renewcommand*{\chaptitlefont}{\normalfont\LARGE\bfseries\color{primary}\raggedright}
  \renewcommand*{\printchaptername}{\chapnamefont \MakeUppercase{\chaptername}}
  \renewcommand*{\printchapternum}{\chapnumfont \thechapter}
  \renewcommand*{\afterchapternum}{\par\nobreak\vskip 5pt\color{boxborder}\hrule height 1.5pt\vskip 15pt}
  \renewcommand*{\printchaptertitle}[1]{\chaptitlefont ##1}
}
\chapterstyle{manualtheme}

% --- Estilo de Cabeçalhos e Rodapés ---
\makepagestyle{manualheadings}
\makeevenhead{manualheadings}{\small\color{secondary}\thepage}{}{\small\color{secondary}\leftmark}
\makeoddhead{manualheadings}{\small\color{secondary}\rightmark}{}{\small\color{secondary}\thepage}
\makeevenfoot{manualheadings}{}{}{}
\makeoddfoot{manualheadings}{}{}{}
\makeheadrule{manualheadings}{\textwidth}{0.4pt}

\pagestyle{manualheadings}

% Formatação das seções e subseções
\setsecheadstyle{\Large\bfseries\color{primary}}
\setsubsecheadstyle{\large\bfseries\color{secondary}}
\setsubsubsecheadstyle{\normalsize\bfseries\color{neutraldark}}

% --- Hiperlinks (carregado antes de abntex2cite) ---
\usepackage{hyperref}
\usepackage{xurl}
\hypersetup{
  colorlinks=true,
  linkcolor=primary,
  citecolor=primary,
  filecolor=secondary,
  urlcolor=secondary,
  bookmarksnumbered=true,
  pdfauthor={Equipe de Forma{\c c}{\~a}o Docente e IA no Ensino Superior},
  pdftitle={Manual de Skills e Agentes para OpenCode: Avalia{\c c}{\~a}o Formativa e Intelig{\^e}ncia Artificial na Educa{\c c}{\~a}o},
  pdfkeywords={Educa{\c c}{\~a}o Superior, Intelig{\^e}ncia Artificial, Avalia{\c c}{\~a}o Formativa, OpenCode, Skills, Agentes, Escala AIAS, Taxonomia de Bloom, Rubricas}
}

% --- Citações no Formato ABNT ---
\usepackage[alf]{abntex2cite}

% O abntex2cite impoe os delimitadores < > ao comando \path (do pacote url),
% exibindo <slug> ao inves de slug. Restauramos os delimitadores vazios.
\def\UrlLeft{}
\def\UrlRight{}

% Ajuste no recuo do parágrafo
\setlength{\parindent}{1.5cm}
\setlength{\parskip}{0.3em}

% =========================================================================
% INÍCIO DO DOCUMENTO
% =========================================================================
\begin{document}

\frontmatter

% --- Capa Elegante ---
\begin{titlingpage}
  \begin{center}
    \vspace*{1.5cm}
    {\large\scshape\bfseries\color{primary} UNIVERSIDADE FEDERAL DE ITAJUB{\'A} -- UNIFEI}\\[0.3cm]
    {\small\scshape\color{secondary} Acervo Documental e Tecnol{\'o}gico de Intelig{\^e}ncia Artificial na Educa{\c c}{\~a}o}\\[2.5cm]

    \color{boxborder}\rule{\linewidth}{2pt}\\[0.5cm]
    {\huge\bfseries\color{primary} Manual de Skills e Agentes\\ para OpenCode}\\[0.4cm]
    {\Large\color{secondary} Integra{\c c}{\~a}o Pedag{\'o}gica, Avalia{\c c}{\~a}o Formativa e\\ Racioc{\'i}nio Aumentado por IA}\\[0.4cm]
    \color{boxborder}\rule{\linewidth}{2pt}\\[3cm]

    \begin{minipage}{0.85\linewidth}
      \centering
      \itshape\color{neutraldark}
      Guia te{\'o}rico e pr{\'a}tico para docentes e desenvolvedores educacionais.
      Aborda fundamentos de avalia{\c c}{\~a}o formativa, Taxonomia de Bloom Revisada,
      rubricas anal{\'i}ticas, a Escala AIAS e a opera{\c c}{\~a}o do CLI OpenCode com o acervo
      de 62 skills tem{\'a}ticas e 6 subagentes especializados.
    \end{minipage}

    \vfill
    {\small\bfseries\color{secondary} Itajub{\'a} -- MG}\\[0.2cm]
    {\small\color{secondary} 2026}
  \end{center}
\end{titlingpage}

% --- Sumário ---
\tableofcontents*
\clearpage

\mainmatter

% --- Inclusão dos Capítulos ---
\chapter{Introdu{\c c}{\~a}o {\`a} Avalia{\c c}{\~a}o Formativa e {\`a} IA}
\label{cap:introducao}

\section{O Prop{\'o}sito deste Manual e do Reposit{\'o}rio}

O reposit{\'o}rio \texttt{skills-ia-educacao} constitui um acervo documental e tecnol{\'o}gico concebido para preencher a lacuna entre as diretrizes normativas (como o Referencial do MEC \cite{brasil2026referencial} e as resolu{\c c}{\~o}es universit{\'a}rias, por exemplo, da UNIFEI \cite{unifei2025graduacao}) e a pr{\'a}tica cotidiana do corpo docente. Trata-se de produto do \textbf{Projeto de Ensino AvalIA -- Kit de Avalia{\c c}{\~a}o Formativa com Intelig{\^e}ncia Artificial para Professores do Ensino Superior}, desenvolvido com o apoio da Pr{\'o}-Reitoria de Gradua{\c c}{\~a}o da UNIFEI: sua finalidade {\'e} instrumentar o corpo docente com recursos prontos de avalia{\c c}{\~a}o formativa mediada por IA, materializando as diretrizes institucionais em ferramentas de uso imediato.

Em vez de exigir que cada professor se torne um especialista em engenharia de \emph{prompts} complexos, o reposit{\'o}rio disponibiliza \textbf{62 skills tem{\'a}ticas}, reutiliz{\'a}veis em CLIs compat{\'i}veis com o padr{\~a}o Agent Skills, e \textbf{6 subagentes} configurados especificamente para operar no \textbf{OpenCode} (\url{https://opencode.ai}).

\begin{conceitochave}
\textbf{OpenCode, skill e subagente em tr{\^e}s frases}: pense no \textbf{OpenCode} como a oficina (o programa que executa as tarefas); na \textbf{skill} como a ficha de instru{\c c}{\~o}es de uma tarefa (ex.: \texttt{/ia-educacao-rubrica} ensina a montar uma rubrica); e no \textbf{subagente} como o especialista que segue v{\'a}rias fichas para entregar um trabalho completo (ex.: \texttt{@construtor-de-avaliacoes} monta prova com rubrica e gabarito). Voc{\^e} n{\~a}o precisa decorar nada: basta descrever o que quer ou digitar a barra (\texttt{/}) para skills e o arroba (\texttt{@}) para subagentes.
\end{conceitochave}

\begin{objetivos}
\begin{itemize}
  \item Compreender por que a avalia{\c c}{\~a}o formativa se torna central quando a IA generativa transforma as tarefas acad{\^e}micas;
  \item Conhecer os cinco n{\'i}veis da Escala AIAS e diferenciar seu car{\'a}ter n{\~a}o hier{\'a}rquico de outras taxonomias;
  \item Visualizar o percurso completo do manual e saber a que cap{\'i}tulo recorrer em cada necessidade docente.
\end{itemize}
\end{objetivos}

\section{Rota R{\'a}pida para Come{\c c}ar}

Se esta {\'e} a sua primeira experi{\^e}ncia com IA generativa, n{\~a}o {\'e} necess{\'a}rio dominar todo o manual antes de realizar a primeira atividade. O percurso abaixo leva cerca de 60 minutos e termina com um material rascunhado, revisado e com a pol{\'i}tica de IA declarada. Siga nesta ordem:

\begin{enumerate}
  \item Leia este cap{\'i}tulo e o Cap{\'i}tulo~\ref{cap:fundamentos} para compreender a finalidade pedag{\'o}gica da IA (cerca de 15 minutos; leitura conceitual, sem computador);
  \item Instale o OpenCode, configure um provedor e instale as skills e os subagentes seguindo o Cap{\'i}tulo~\ref{cap:instalacao} (cerca de 20 minutos; use apenas o \emph{caminho recomendado} do seu sistema e ignore os m{\'e}todos alternativos por enquanto);
  \item Execute um caso simples do Cap{\'i}tulo~\ref{cap:casos-de-uso}, substituindo a disciplina e o tema pelos seus (cerca de 15 minutos; o Caso~2 {\'e} uma boa primeira escolha para quem avalia, o Caso~1 para quem planeja);
  \item Revise o resultado com o checklist do pr{\'o}prio Cap{\'i}tulo~\ref{cap:casos-de-uso} antes de entregar qualquer material aos estudantes (cerca de 5 minutos; nada vai aos estudantes sem passar por esta revis{\~a}o);
  \item Use o Cap{\'i}tulo~\ref{cap:explicar-estudantes} para declarar a pol{\'i}tica de uso de IA da turma (cerca de 5 minutos; copie o texto pronto de declara{\c c}{\~a}o). Os Cap{\'i}tulos~\ref{cap:skills-opencode} e~\ref{cap:agentes-skills} s{\~a}o refer{\^e}ncia de consulta: podem ficar para depois, quando voc{\^e} quiser entender ou manter o acervo.
\end{enumerate}

\subsection{Vocabul{\'a}rio M{\'i}nimo}

\begin{itemize}
  \item \textbf{IAGen}: intelig{\^e}ncia artificial generativa, capaz de produzir texto, c{\'o}digo, imagens ou outros artefatos a partir de instru{\c c}{\~o}es;
  \item \textbf{LLM}: modelo de linguagem de grande escala, usado para prever e gerar sequ{\^e}ncias de texto;
  \item \textbf{Prompt}: instru{\c c}{\~a}o, contexto e crit{\'e}rios fornecidos ao modelo;
  \item \textbf{Skill}: arquivo de instru{\c c}{\~o}es especializadas carregado sob demanda;
  \item \textbf{Subagente}: especialista configurado para coordenar uma tarefa complexa e, conforme suas permiss{\~o}es, usar v{\'a}rias skills;
  \item \textbf{AIAS}: escala que explicita como a IA pode participar de uma atividade avaliativa;
  \item \textbf{STHEM}: sigla para as grandes {\'a}reas Science, Technology, Humanities, Engineering e Mathematics (Ci{\^e}ncias, Tecnologia, Humanidades, Engenharias e Matem{\'a}tica);
  \item \textbf{DUA}: abordagem que reduz barreiras oferecendo m{\'u}ltiplos meios de engajamento, representa{\c c}{\~a}o e express{\~a}o;
  \item \textbf{CoVe}: procedimento de verifica{\c c}{\~a}o independente de afirma{\c c}{\~o}es produzidas por IA;
  \item \textbf{CLI e TUI}: \emph{CLI} {\'e} o terminal (tela escura de comandos); \emph{TUI} {\'e} a interface interativa que abre ao digitar \texttt{opencode} no terminal;
  \item \textbf{Provedor e modelo}: o \emph{provedor} {\'e} a empresa ou servi{\c c}o que oferece a IA; o \emph{modelo} {\'e} o ``c{\'e}rebro'' escolhido na hora (ver Cap{\'i}tulo~\ref{cap:modelos-gratuitos});
  \item \textbf{Headless}: modo de usar o OpenCode com uma {\'u}nica linha de comando, sem abrir a interface (ex.: \texttt{opencode run "..."});
  \item \textbf{Frontmatter}: pequeno cabe{\c c}alho t{\'e}cnico no in{\'i}cio de cada skill; voc{\^e} n{\~a}o precisa edit{\'a}-lo para usar as skills;
  \item \textbf{MCP}: padr{\~a}o aberto que liga o OpenCode a ferramentas externas; citado aqui s{\'o} para registro de proced{\^e}ncia;
  \item \textbf{ConcepTest e blueprint}: \emph{ConcepTest} {\'e} uma quest{\~a}o conceitual curta para vota{\c c}{\~a}o em sala; \emph{blueprint} {\'e} o mapa da prova (quantas quest{\~o}es por t{\'o}pico e n{\'i}vel).
\end{itemize}

Ao longo dos pr{\'o}ximos cap{\'i}tulos, este manual apresentar{\'a} em detalhes:
\begin{itemize}
  \item Os fundamentos pedag{\'o}gicos de feedback, taxonomia cognitiva e rubricas anal{\'i}ticas (Cap{\'i}tulo~\ref{cap:fundamentos});
  \item A arquitetura das skills e seu mecanismo de execu{\c c}{\~a}o com OpenCode (Cap{\'i}tulo~\ref{cap:skills-opencode});
  \item O passo a passo completo de instala{\c c}{\~a}o do OpenCode e sincroniza{\c c}{\~a}o do acervo (Cap{\'i}tulo~\ref{cap:instalacao});
  \item O fluxo de trabalho na interface gr{\'a}fica e desktop (Cap{\'i}tulo~\ref{cap:desktop});
  \item A escolha de modelos gratuitos no OpenCode (Cap{\'i}tulo~\ref{cap:modelos-gratuitos});
  \item Princ{\'i}pios de {\'e}tica, limita{\c c}{\~o}es, LGPD e transpar{\^e}ncia (Cap{\'i}tulo~\ref{cap:etica-transparencia});
  \item Estrat{\'e}gias para explicar a IA e a Escala AIAS aos estudantes (Cap{\'i}tulo~\ref{cap:explicar-estudantes});
  \item A especifica{\c c}{\~a}o dos seis subagentes e o cat{\'a}logo das 62 skills (Cap{\'i}tulo~\ref{cap:agentes-skills});
  \item Casos de uso iniciais e roteiros pr{\'a}ticos para sala de aula (Cap{\'i}tulo~\ref{cap:casos-de-uso}).
\end{itemize}

\section{O Desafio Pedag{\'o}gico na Era da Intelig{\^e}ncia Artificial}

A r{\'a}pida difus{\~a}o de modelos de linguagem de grande escala (\emph{Large Language Models} -- LLMs) e de ferramentas baseadas em Intelig{\^e}ncia Artificial Generativa (IAGen) vem transformando a educa{\c c}{\~a}o b{\'a}sica e superior. Tarefas cl{\'a}ssicas de avalia{\c c}{\~a}o, como reda{\c c}{\~a}o de ensaios curtos, resolu{\c c}{\~a}o de quest{\~o}es conceituais elementares e gera{\c c}{\~a}o de trechos introdut{\'o}rios de c{\'o}digo-fonte, podem ser realizadas com facilidade por assistentes computacionais \cite{unesco2021guidance,mollick2024instructors}.

Frente a esse cen{\'a}rio, rea{\c c}{\~o}es puramente proibitivas ou tentativas ing{\^e}nuas de detec{\c c}{\~a}o automatizada de conte{\'u}do gerado por IA demonstraram-se ineficazes, eticamente arriscadas e pedagogicamente est{\'e}reis. Como alerta o Minist{\'e}rio da Educa{\c c}{\~a}o em seu documento norteador, o \emph{Referencial para Desenvolvimento e Uso Respons{\'a}veis de Intelig{\^e}ncia Artificial na Educa{\c c}{\~a}o} \cite{brasil2026referencial}, a intelig{\^e}ncia artificial deve ser compreendida n{\~a}o como uma amea{\c c}a a ser banida, mas como um catalisador de revis{\~a}o das pr{\'a}ticas docentes, exigindo centralidade humana, transpar{\^e}ncia {\'e}tica e intencionalidade pedag{\'o}gica.

A superada {\^e}nfase em produtos finais padronizados precisa ceder espa{\c c}o para a valoriza{\c c}{\~a}o dos processos de pensamento, da investiga{\c c}{\~a}o cr{\'i}tica, da metacogni{\c c}{\~a}o e da avalia{\c c}{\~a}o formativa cont{\'i}nua.

\section{Da Avalia{\c c}{\~a}o Somativa {\`a} Avalia{\c c}{\~a}o Formativa}

Historicamente, as institui{\c c}{\~o}es de ensino superior priorizaram a \textbf{avalia{\c c}{\~a}o somativa} -- instrumentos pontuais aplicados ao t{\'e}rmino de uma unidade ou semestre (provas tradicionais, exames finais e testes de m{\'u}ltipla escolha sem devolutiva) cujo objetivo prec{\'i}puo {\'e} certificar, classificar ou atribuir uma nota ao estudante. 

Embora a certifica{\c c}{\~a}o cumpra papel institucional regulat{\'o}rio, a literatura cient{\'i}fica internacional h{\'a} d{\'e}cadas demonstra que a avalia{\c c}{\~a}o somativa isolada possui impacto limitado sobre a qualidade efetiva da aprendizagem \cite{black1998assessment}. Em contrapartida, a \textbf{avalia{\c c}{\~a}o formativa} concebe a avalia{\c c}{\~a}o como parte intr{\'i}nseca do ato de ensinar e aprender. 

\begin{conceitochave}
\textbf{Avalia{\c c}{\~a}o Formativa}: Qualquer processo avaliativo planejado e implementado ao longo do percurso de aprendizagem com a finalidade de coletar evid{\^e}ncias sobre a compreens{\~a}o dos estudantes, fornecendo informa{\c c}{\~o}es que permitam a professores e alunos ajustarem o ensino e as estrat{\'e}gias de estudo em tempo real \cite{black1998assessment}.
\end{conceitochave}

Segundo a formula{\c c}{\~a}o seminal de \citeonline{sadler1989formative}, a avalia{\c c}{\~a}o formativa eficaz exige que o estudante compreenda tr{\^e}s elementos fundamentais:
\begin{enumerate}
  \item \textbf{O padr{\~a}o ou objetivo de refer{\^e}ncia}: clareza sobre o que constitui um desempenho de excel{\^e}ncia;
  \item \textbf{O estado atual de aprendizagem}: reconhecimento cr{\'i}tico de onde o pr{\'o}prio trabalho se situa no momento;
  \item \textbf{As a{\c c}{\~o}es de fechamento do hiato (\emph{closing the gap})}: estrat{\'e}gias claras e acion{\'a}veis para mover o desempenho atual em dire{\c c}{\~a}o ao padr{\~a}o desejado.
\end{enumerate}

Em ambientes educacionais permeados por IA, essa concep{\c c}{\~a}o ganha urg{\^e}ncia. Se uma m{\'a}quina pode redigir a resposta final, o valor pedag{\'o}gico desloca-se para a capacidade do aluno em interrogar o problema, planejar abordagens, validar solu{\c c}{\~o}es, corrigir inconsist{\^e}ncias e justificar decis{\~o}es tomadas.

\section{A Escala AIAS: Transpar{\^e}ncia e Intencionalidade Pedag{\'o}gica}

Para que a integra{\c c}{\~a}o da intelig{\^e}ncia artificial n{\~a}o ocorra em um v{\'a}cuo normativo ou sob regras amb{\'i}guas, a literatura contempor{\^a}nea desenvolveu a \textbf{AIAS} (\emph{Artificial Intelligence Assessment Scale}), proposta originalmente por \citeonline{perkins2024aias} e refinada por \citeonline{perkins2025reimagining}.

A AIAS {\'e} um referencial conceitual que classifica o n{\'i}vel de uso de IA admitido ou exigido em uma determinada atividade avaliativa. Ela n{\~a}o se destina a vigiar o aluno, mas sim a explicitar o contrato pedag{\'o}gico entre docente e discente. A Tabela~\ref{tab:aias-escala} apresenta os cinco n{\'i}veis can{\^o}nicos adotados estritamente em todo o ecossistema deste reposit{\'o}rio.

\begin{table}[htbp]
  \centering
  \small
  \caption{Os 5 N{\'i}veis Can{\^o}nicos da Escala AIAS}
  \label{tab:aias-escala}
  \begin{tabularx}{\linewidth}{clp{4cm}X}
    \toprule
    \textbf{N{\'i}vel} & \textbf{Nome Can{\^o}nico} & \textbf{IA Permitida} & \textbf{Produto Final Esperado} \\
    \midrule
    \textbf{1} & Sem IA & Nenhuma interven{\c c}{\~a}o computacional inteligente & Elaborado 100\% pelo estudante sem qualquer aux{\'i}lio de IA. \\
    \addlinespace
    \textbf{2} & Planejamento Assistido por IA & Idea{\c c}{\~a}o, \emph{brainstorming}, esbo{\c c}os e estrutura{\c c}{\~a}o & Do estudante; a IA atua apenas na fase inicial preparat{\'o}ria. \\
    \addlinespace
    \textbf{3} & Colabora{\c c}{\~a}o com IA & Elabora{\c c}{\~a}o conjunta, refinamento e expans{\~a}o & Do estudante com aux{\'i}lio expl{\'i}cito de IA em trechos espec{\'i}ficos. \\
    \addlinespace
    \textbf{4} & IA Integral & Uso estrat{\'e}gico, abrangente e cont{\'i}nuo & Dirigido criticamente pelo estudante com forte presen{\c c}a de IA. \\
    \addlinespace
    \textbf{5} & Explora{\c c}{\~a}o de IA & Co-cria{\c c}{\~a}o avan{\c c}ada e inova{\c c}{\~a}o de fronteira & Parceria autoral entre estudante e IA em problemas complexos. \\
    \bottomrule
  \end{tabularx}
\end{table}

Duas propriedades estruturais da Escala AIAS s{\~a}o essenciais e devem orientar todo o desenho instrucional:
\begin{itemize}
  \item \textbf{Car{\'a}ter n{\~a}o hier{\'a}rquico}: Nenhum n{\'i}vel {\'e} intrinsecamente superior ou mais nobre do que outro. Uma atividade no N{\'i}vel~1 (Sem IA) pode demandar criatividade extrema e racioc{\'i}nio anal{\'i}tico profundo (por exemplo, um debate oral ou uma solu{\c c}{\~a}o manuscrita em sala), enquanto uma atividade no N{\'i}vel~4 pode focar em habilidades de auditoria cr{\'i}tica de c{\'o}digo ou revis{\~a}o de relat{\'o}rios. A escolha do n{\'i}vel deriva estritamente dos \textbf{objetivos de aprendizagem} visados;
  \item \textbf{Car{\'a}ter cumulativo}: N{\'i}veis superiores autorizam as pr{\'a}ticas dos n{\'i}veis inferiores, salvo proibi{\c c}{\~a}o expressa no enunciado. Por exemplo, uma atividade AIAS~3 admite que o aluno realize idea{\c c}{\~a}o inicial com a ferramenta (pr{\'a}tica do N{\'i}vel~2), avan{\c c}ando para o co-refinamento de se{\c c}{\~o}es espec{\'i}ficas.
\end{itemize}

\begin{atencaopedagogica}
\textbf{Declara{\c c}{\~a}o Pr{\'e}via Obrigat{\'o}ria}: O n{\'i}vel AIAS deve ser declarado de forma inequ{\'i}voca no enunciado da atividade e no plano de ensino \emph{antes} do in{\'i}cio da tarefa. Exigir do aluno uma postura em rela{\c c}{\~a}o {\`a} IA sem defini{\c c}{\~a}o formal pr{\'e}via compromete a seguran{\c c}a jur{\'i}dica e pedag{\'o}gica da avalia{\c c}{\~a}o.
\end{atencaopedagogica}

Ao planejar uma atividade, separe duas perguntas: a IA generativa est{\'a} \textbf{permitida} e a IA generativa est{\'a} \textbf{obrigat{\'o}ria}? O n{\'i}vel AIAS informa o papel da IA no processo, mas o enunciado deve esclarecer tamb{\'e}m quais evid{\^e}ncias ser{\~a}o avaliadas, como o uso ser{\'a} declarado e quais tecnologias assistivas ou adapta{\c c}{\~o}es institucionais permanecem autorizadas. ``Sem IA'' deve ser entendido como sem IA generativa na tarefa, n{\~a}o como proibi{\c c}{\~a}o de acessibilidade.

\chapter{Fundamentos Pedag{\'o}gicos}
\label{cap:fundamentos}

\begin{objetivos}
\begin{itemize}
  \item Aplicar o modelo tridimensional de feedback (Feed Up, Feed Back, Feed Forward) na devolutiva de trabalhos acad{\^e}micos;
  \item Classificar objetivos de aprendizagem nos seis n{\'i}veis da Taxonomia Revisada de Bloom;
  \item Estruturar rubricas anal{\'i}ticas com r{\'o}tulos desenvolvimentais, incluindo o crit{\'e}rio de transpar{\^e}ncia no uso de IA;
  \item Relacionar o planejamento reverso (Backward Design) e o Desenho Universal para a Aprendizagem ao desenho de disciplinas.
\end{itemize}
\end{objetivos}

\section{O Feedback Formativo: Conceito e Mecanismos de A{\c c}{\~a}o}

O feedback {\'e} reconhecido pela literatura em educa{\c c}{\~a}o como uma das interven{\c c}{\~o}es pedag{\'o}gicas com maior pot{\^e}ncia para impulsionar o desempenho discente \cite{hattie2007power}. Todavia, a efici{\^e}ncia do feedback n{\~a}o {\'e} uniforme: h{\'a} modalidades que aceleram a aprendizagem, enquanto interven{\c c}{\~o}es mal estruturadas podem gerar desengajamento, ansiedade ou paralisia.

\subsection{O Modelo Tridimensional de Hattie \& Timperley}

Em seu estudo seminal, \citeonline{hattie2007power} estabeleceram que o feedback verdadeiramente formativo precisa responder, de modo cont{\'i}nuo e integrado, a tr{\^e}s quest{\~o}es centrais:

\begin{enumerate}
  \item \textbf{\emph{Feed Up} (Para onde vou?)}: Refere-se {\`a} clareza sobre as metas de aprendizagem, os crit{\'e}rios de sucesso e o padr{\~a}o de desempenho esperado. Se o estudante n{\~a}o sabe exatamente o que precisa demonstrar, qualquer devolu{\c c}{\~a}o torna-se incompreens{\'i}vel;
  \item \textbf{\emph{Feed Back} (Como estou indo?)}: Consiste na an{\'a}lise descritiva e objetiva da produ{\c c}{\~a}o atual do estudante em confronto direto com os crit{\'e}rios estabelecidos, evidenciando o que j{\'a} foi atingido e o que apresenta desconformidade;
  \item \textbf{\emph{Feed Forward} (O que fazer a seguir?)}: Constitui o componente mais frequentemente negligenciado nas salas de aula tradicionais. Representa o roteiro de a{\c c}{\~o}es imediatas, concretas e pratic{\'a}veis que o estudante deve executar para fechar a lacuna identificada e avan{\c c}ar em complexidade.
\end{enumerate}

\subsection{Os Quatro N{\'i}veis de Foco do Feedback}

Al{\'e}m das tr{\^e}s quest{\~o}es, \citeonline{hattie2007power} categorizam o feedback em quatro n{\'i}veis hier{\'a}rquicos de foco:

\begin{itemize}
  \item \textbf{Foco na Tarefa (FT)}: Informa se a resposta est{\'a} correta ou incorreta e aponta informa{\c c}{\~o}es factuais faltantes. {\'E} {\'u}til para iniciantes, mas limitado para transfer{\^e}ncia de conhecimento;
  \item \textbf{Foco no Processo (FP)}: Examina as estrat{\'e}gias cognitivas, os m{\'e}todos de resolu{\c c}{\~a}o, as rela{\c c}{\~o}es causais e os passos l{\'o}gicos adotados pelo estudante. Apresenta alta efic{\'a}cia para o desenvolvimento de racioc{\'i}nio profundo;
  \item \textbf{Foco na Autorregula{\c c}{\~a}o (FR)}: Estimula o monitoramento aut{\^o}nomo, a autocorre{\c c}{\~a}o, a gest{\~a}o do tempo e a confian{\c c}a reflexiva do aluno. Apresenta o mais duradouro impacto na forma{\c c}{\~a}o acad{\^e}mica e profissional;
  \item \textbf{Foco no \emph{Self} / Pessoa (FS)}: Consiste em elogios ou cr{\'i}ticas gen{\'e}ricas dirigidas {\`a} pessoa do estudante (ex.: ``Voc{\^e} {\'e} muito inteligente'' ou ``Voc{\^e} n{\~a}o tem voca{\c c}{\~a}o para c{\'a}lculo''). A literatura evidencia que o feedback no n{\'i}vel da pessoa possui impacto pedag{\'o}gico nulo ou negativo, pois desvia a aten{\c c}{\~a}o da tarefa e refor{\c c}a mentalidades fixas de intelig{\^e}ncia.
\end{itemize}

\begin{conceitochave}
\textbf{Princ{\'i}pio da Carga Cognitiva no Feedback}: Um feedback eficaz limita-se a um n{\'u}mero gerenci{\'a}vel de a{\c c}{\~o}es de melhoria (idealmente entre uma e duas a{\c c}{\~o}es priorit{\'a}rias). Relat{\'o}rios com dezenas de corre{\c c}{\~o}es simult{\^a}neas sobrecarregam a mem{\'o}ria de trabalho do aluno e resultam em abandono.
\end{conceitochave}

Complementarmente, \citeonline{nicol2006formative} salientam que o bom feedback deve estimular o di{\'a}logo e a autoavalia{\c c}{\~a}o cr{\'i}tica discente. Em nossa arquitetura, os agentes e as skills de feedback operam com base estrita nesses princ{\'i}pios: analisam processos, formulam perguntas ativadoras da metacogni{\c c}{\~a}o e estruturam passos de \emph{feed forward} claros e vi{\'a}veis.

\section{A Taxonomia Revisada de Bloom}

A classifica{\c c}{\~a}o de objetivos educacionais proposta originalmente por \citeonline{bloom1956taxonomy} foi substancialmente reformulada por \citeonline{anderson2001taxonomy}. A \textbf{Taxonomia Revisada de Bloom (TRB)} converteu os substantivos originais em formas verbais din{\^a}micas e estruturou o dom{\'i}nio cognitivo em uma matriz bidimensional: a \emph{Dimens{\~a}o do Processo Cognitivo} e a \emph{Dimens{\~a}o do Conhecimento}.

\subsection{A Dimens{\~a}o do Processo Cognitivo}

A progress{\~a}o cognitiva compreende seis n{\'i}veis em ordem crescente de complexidade, conforme sintetizado na Tabela~\ref{tab:bloom-niveis}.

\begin{table}[htbp]
  \centering
  \small
  \caption{Os 6 N{\'i}veis Cognitivos da Taxonomia Revisada de Bloom e o Impacto da IA}
  \label{tab:bloom-niveis}
  \begin{tabularx}{\linewidth}{clp{4.5cm}X}
    \toprule
    \textbf{N{\'i}vel} & \textbf{Processo} & \textbf{Verbos de A{\c c}{\~a}o} & \textbf{Vulnerabilidade / Oportunidade com IA} \\
    \midrule
    \textbf{1} & Lembrar & Listar, reconhecer, definir, nomear, reproduzir & \emph{Alta vulnerabilidade}. A IA responde instantaneamente; valor pedag{\'o}gico isolado reduzido. \\
    \addlinespace
    \textbf{2} & Entender & Explicar, exemplificar, resumir, interpretar, classificar & \emph{M{\'e}dia vulnerabilidade}. LLMs parafraseiam conceitos; exige exig{\^e}ncia de exemplos pr{\'o}prios e locais. \\
    \addlinespace
    \textbf{3} & Aplicar & Implementar, resolver, executar, calcular, usar & \emph{M{\'e}dia-alta}. A IA executa procedimentos padr{\~a}o; requer problemas contextualizados e dados novos. \\
    \addlinespace
    \textbf{4} & Analisar & Comparar, diferenciar, organizar, desconstruir, auditar & \emph{Alta oportunidade}. O aluno analisa erros em c{\'o}digos ou relat{\'o}rios gerados pela pr{\'o}pria IA. \\
    \addlinespace
    \textbf{5} & Avaliar & Julgar, criticar, defender, priorizar, validar & \emph{Elevada relev{\^a}ncia}. Demanda julgamento de valor, an{\'a}lise de vieses {\'e}ticos e conformidade normativa. \\
    \addlinespace
    \textbf{6} & Criar & Projetar, formular, planejar, sintetizar, construir & \emph{Parceria criativa}. Co-cria{\c c}{\~a}o de solu{\c c}{\~o}es originais onde o estudante assume a dire{\c c}{\~a}o autoral. \\
    \bottomrule
  \end{tabularx}
\end{table}

\subsection{A Dimens{\~a}o do Conhecimento}

Paralelamente aos processos mentais, os objetivos educacionais articulam quatro tipos de conhecimento \cite{anderson2001taxonomy}:
\begin{enumerate}
  \item \textbf{Conhecimento Factual}: Elementos b{\'a}sicos, terminologias e detalhes espec{\'i}ficos;
  \item \textbf{Conhecimento Conceitual}: Esquemas, categorias, princ{\'i}pios e generaliza{\c c}{\~o}es estruturantes;
  \item \textbf{Conhecimento Procedimental}: M{\'e}todos, t{\'e}cnicas, algoritmos e crit{\'e}rios para aplicar solu{\c c}{\~o}es;
  \item \textbf{Conhecimento Metacognitivo}: Consci{\^e}ncia da pr{\'o}pria cogni{\c c}{\~a}o, conhecimento estrat{\'e}gico e gest{\~a}o do pr{\'o}prio aprendizado.
\end{enumerate}

\begin{conceitochave}
\textbf{Independ{\^e}ncia entre Bloom e AIAS}: {\'E} fundamental n{\~a}o confundir o n{\'i}vel de Bloom com o n{\'i}vel AIAS. Uma atividade no n{\'i}vel AIAS~1 (Sem IA) pode demandar processos de ordem superior (\emph{Criar} ou \emph{Avaliar}), como uma prova pr{\'a}tica de bancada ou reda{\c c}{\~a}o reflexiva manuscrita. Inversamente, uma atividade AIAS~4 (IA Integral) pode focar no n{\'i}vel cognitivo \emph{Analisar}, exigindo que o aluno audite e encontre vulnerabilidades de seguran{\c c}a em um c{\'o}digo gerado por IA.
\end{conceitochave}

\section{Rubricas de Avalia{\c c}{\~a}o Formativa}

Uma rubrica {\'e} um guia pontuado e estruturado que combina crit{\'e}rios de avalia{\c c}{\~a}o com descritores qualitativos de desempenho para diferentes n{\'i}veis de dom{\'i}nio \cite{brookhart2013rubrics}. As rubricas cumprem papel decisivo na avalia{\c c}{\~a}o formativa:
\begin{itemize}
  \item Comunicam previamente as expectativas aos estudantes;
  \item Reduzem a subjetividade e o vi{\'e}s inconsciente na corre{\c c}{\~a}o;
  \item Possibilitam a autoavalia{\c c}{\~a}o e a avalia{\c c}{\~a}o por pares com rigor \cite{andrade2019critical};
  \item Transformam a corre{\c c}{\~a}o em um mapa claro de \emph{feed back} e \emph{feed forward}.
\end{itemize}

\subsection{Estrutura de uma Rubrica Anal{\'i}tica Robusta}

Conforme prescreve a skill \texttt{ia-educacao-rubrica}, as rubricas desenvolvidas no reposit{\'o}rio s{\~a}o prioritariamente \textbf{anal{\'i}ticas} (avaliam crit{\'e}rios individualmente) e seguem crit{\'e}rios rigorosos de desenho:

\begin{enumerate}
  \item \textbf{Espelhamento de Objetivos}: Cada crit{\'e}rio da rubrica deve espelhar diretamente um objetivo de aprendizagem declarado. Crit{\'e}rios cosm{\'e}ticos desprovidos de respaldo formativo (ex.: ``est{\'e}tica geral'') s{\~a}o desencorajados;
  \item \textbf{Extens{\~a}o Gerenci{\'a}vel}: Recomenda-se de 3 a 6 crit{\'e}rios por rubrica, com 4 n{\'i}veis de desempenho;
  \item \textbf{R{\'o}tulos Neutros e Desenvolvimentais}: Adotam-se r{\'o}tulos que valorizam o percurso de aprendizagem:
    \begin{itemize}
      \item \emph{Iniciante}: O estudante compreende os objetivos b{\'a}sicos, mas apresenta inconsist{\^e}ncias conceituais ou lacunas t{\'e}cnicas substanciais;
      \item \emph{Em desenvolvimento}: Demonstra compreens{\~a}o parcial e aplica procedimentos essenciais, necessitando de ajustes e refinamentos;
      \item \emph{Proficiente}: Atende plenamente aos objetivos de aprendizagem com rigor t{\'e}cnico, clareza e autonomia;
      \item \emph{Exemplar}: Supera as expectativas m{\'i}nimas, integrando an{\'a}lise cr{\'i}tica original, conex{\~o}es interdisciplinares e excel{\^e}ncia metodol{\'o}gica.
    \end{itemize}
  \item \textbf{Crit{\'e}rio de Coer{\^e}ncia com a Escala AIAS}: Em tarefas que admitem ou exigem ferramentas computacionais, a rubrica inclui um crit{\'e}rio espec{\'i}fico sobre transpar{\^e}ncia e integridade no uso de IA, avaliando se a autoria e o refinamento respeitaram o n{\'i}vel AIAS acordado.
\end{enumerate}

\section{Planejamento Did{\'a}tico e Inclus{\~a}o}

\subsection{Planejamento Reverso (\emph{Backward Design})}

A metodologia de \emph{Backward Design}, formulada por \citeonline{wiggins2005understanding}, estabelece que o planejamento did{\'a}tico deve come{\c c}ar n{\~a}o pela escolha de conte{\'u}dos ou listas de livros, mas pelo destino final da aprendizagem. A estrutura organiza-se em tr{\^e}s etapas sequenciais:

\begin{enumerate}
  \item \textbf{Etapa 1 -- Identificar os Resultados Desejados}: O que os estudantes devem ser capazes de saber, compreender e fazer ao t{\'e}rmino do percurso? Quais s{\~a}o as ideias perenes e compreens{\~o}es essenciais?
  \item \textbf{Etapa 2 -- Determinar as Evid{\^e}ncias Aceit{\'a}veis}: Como saberemos se os estudantes alcan{\c c}aram esses resultados? Quais tarefas aut{\^e}nticas, avalia{\c c}{\~o}es formativas e rubricas fornecer{\~a}o evid{\^e}ncias leg{\'i}timas de aprendizagem?
  \item \textbf{Etapa 3 -- Planejar as Experi{\^e}ncias de Aprendizagem e o Ensino}: Quais atividades ativas, leituras e interven{\c c}{\~o}es pedag{\'o}gicas conduzir{\~a}o os discentes {\`a} constru{\c c}{\~a}o das compet{\^e}ncias demandadas?
\end{enumerate}

Esse fluxo dialoga diretamente com o princ{\'i}pio do \textbf{Alinhamento Construtivo} de \citeonline{biggs2011teaching}, segundo o qual o processo de ensino apenas tem {\^e}xito quando as tarefas avaliativas e as estrat{\'e}gias did{\'a}ticas operam em perfeita sintonia cognitiva com os objetivos almejados.

\subsection{Desenho Universal para a Aprendizagem (DUA) e Acessibilidade}

Nenhum sistema educacional pode ser considerado de qualidade se n{\~a}o for inclusivo. O referencial do \textbf{Desenho Universal para a Aprendizagem (DUA)} \cite{cast2018udl} prop{\~o}e a elimina{\c c}{\~a}o sistem{\'a}tica de barreiras pedag{\'o}gicas por meio de tr{\^e}s eixos fundantes:

\begin{itemize}
  \item \textbf{M{\'u}ltiplos Meios de Engajamento} (o ``porqu{\^e}'' da aprendizagem): estimular o interesse, a persist{\^e}ncia e a autorregula{\c c}{\~a}o emocional;
  \item \textbf{M{\'u}ltiplos Meios de Representa{\c c}{\~a}o} (o ``o qu{\^e}'' da aprendizagem): apresentar conceitos por diferentes linguagens (textual, visual, gr{\'a}fica, auditiva e interativa);
  \item \textbf{M{\'u}ltiplos Meios de A{\c c}{\~a}o e Express{\~a}o} (o ``como'' da aprendizagem): oportunizar diferentes formas para que os estudantes demonstrem suas compet{\^e}ncias.
\end{itemize}

No {\^a}mbito de nosso acervo de agentes, a skill \texttt{ia-educacao-dua} e o agente \texttt{coach-de-feedback-formativo} aplicam princ{\'i}pios do DUA na produ{\c c}{\~a}o de feedbacks com linguagem concisa, organiza{\c c}{\~a}o em etapas e redu{\c c}{\~a}o de jarg{\~o}es desnecess{\'a}rios. As adapta{\c c}{\~o}es devem responder {\`a}s necessidades e prefer{\^e}ncias informadas pelo estudante, sem inferir diagn{\'o}sticos. A estrutura das skills (padr{\~a}o \texttt{SKILL.md}) e o cat{\'a}logo completo de subagentes e skills s{\~a}o apresentados nos Cap{\'i}tulos~\ref{cap:skills-opencode} e~\ref{cap:agentes-skills}.

\chapter{O que s{\~a}o Skills? Como Funcionam com o OpenCode?}
\label{cap:skills-opencode}

\section{O Paradigma de Skills em Assistentes de Intelig{\^e}ncia Artificial}

Nas primeiras gera{\c c}{\~o}es de ferramentas baseadas em intelig{\^e}ncia artificial generativa, a intera{\c c}{\~a}o com modelos de linguagem restringia-se ao envio de \emph{prompts} soltos na interface web. Essa abordagem apresentava graves fragilidades para o uso docente e acad{\^e}mico: aus{\^e}ncia de reprodutibilidade, variabilidade excessiva nas respostas, perda constante de diretrizes metodol{\'o}gicas e risco acentuado de alucina{\c c}{\~a}o factual.

Com o amadurecimento dos assistentes de c{\'o}digo e agentes aut{\^o}nomos de terminal -- como o \textbf{OpenCode} (\url{https://opencode.ai}) -- consolidou-se o paradigma de \textbf{Skills} (\emph{habilidades modulares}). 

\begin{conceitochave}
\textbf{Defini{\c c}{\~a}o de Skill}: Uma \emph{skill} {\'e} uma unidade modular, autocontida, versionada e semanticamente estruturada de conhecimento procedimental e diretrizes operacionais. Em vez de injetar previamente manuais enciclop{\'e}dicos no contexto fixo do agente, a skill {\'e} carregada de forma din{\^a}mica e sob demanda (\emph{just-in-time contextual injection}) no exato instante em que o prop{\'o}sito da tarefa discente ou docente {\'e} identificado.
\end{conceitochave}

Essa arquitetura preserva a mem{\'o}ria de trabalho do modelo, reduz custos de infer{\^e}ncia e garante que o comportamento do agente permane{\c c}a rigorosamente balizado por referenciais pedag{\'o}gicos consagrados.

\section{Anatomia de uma Skill: O Padr{\~a}o \texttt{SKILL.md}}

No acervo \texttt{skills-ia-educacao}, cada skill reside em um diret{\'o}rio exclusivo dentro de \texttt{skills/<slug>/} e materializa-se em um documento unificado denominado \texttt{SKILL.md}. A integridade e consist{\^e}ncia de todo o ecossistema repousam sobre duas partes indispens{\'a}veis: o \textbf{Frontmatter YAML} e as \textbf{Oito Se{\c c}{\~o}es Fixas Can{\^o}nicas}.

\subsection{O Frontmatter YAML}

O in{\'i}cio de cada arquivo \texttt{SKILL.md} cont{\'e}m metadados essenciais delimitados por \texttt{---}, conforme exemplificado a seguir:

\begin{lstlisting}[language=bash, caption={Exemplo de Frontmatter YAML de uma Skill}]
---
name: ia-educacao-rubrica
category: ferramentas-praticas
model: any
version: 1.4
description: Especialista em design e construcao de rubricas de avaliacao para ensino superior. Alinha criterios com objetivos de aprendizagem e Taxonomia de Bloom.
---
\end{lstlisting}

Os campos possuem especifica{\c c}{\~o}es r{\'i}gidas auditadas por script:
\begin{itemize}
  \item \texttt{name}: Identificador {\'u}nico (slug) em min{\'u}sculas. Segue o prefixo padronizado \texttt{ia-educacao-<tema>} (com a {\'u}nica exce{\c c}{\~a}o de \texttt{aias-consultant}). O nome do arquivo deve ser rigorosamente id{\^e}ntico ao nome da pasta que o abriga;
  \item \texttt{category}: Classifica{\c c}{\~a}o temática em uma das cinco categorias oficiais do acervo: \texttt{niveis-ensino}, \texttt{formacao-docente}, \texttt{etica-governanca}, \texttt{inclusao-equidade} ou \texttt{ferramentas-praticas};
  \item \texttt{model}: Sempre fixado em \texttt{any}. Garante total independ{\^e}ncia tecnol{\'o}gica, indicando que a l{\'o}gica instrucional funciona com modelos de qualquer provedor (OpenAI, Anthropic, Google, DeepSeek ou modelos locais de c{\'o}digo aberto);
  \item \texttt{version}: Controle sem{\^a}ntico de vers{\~a}o em dois d{\'i}gitos (\texttt{X.Y}), sendo incrementado a cada altera{\c c}{\~a}o metodol{\'o}gica ou estrutural;
  \item \texttt{description}: Texto descritivo conciso e objetivo contendo as palavras-chave e gatilhos que orientam o OpenCode na escolha da skill correta.
\end{itemize}

\subsection{As Oito Se{\c c}{\~o}es Fixas Can{\^o}nicas}

Ap{\'o}s o frontmatter, o corpo do texto {\'e} organizado em exatamente oito se{\c c}{\~o}es Markdown, dispostas obrigatoriamente na seguinte ordem:

\begin{enumerate}
  \item \textbf{\texttt{\#\# Princ{\'i}pios}}: Estabelece a fundamenta{\c c}{\~a}o te{\'o}rico-pedag{\'o}gica e o compromisso {\'e}tico que orientam a a{\c c}{\~a}o docente e a atua{\c c}{\~a}o da IA naquele dom{\'i}nio;
  \item \textbf{\texttt{\#\# Quando usar}}: Define as situa{\c c}{\~o}es t{\'i}picas, gatilhos de contexto e solicita{\c c}{\~o}es do usu{\'a}rio que demandam o acionamento da skill;
  \item \textbf{\texttt{\#\# Workflow}}: Apresenta o passo a passo cognitivo numerado (decomposto em etapas claras) que a IA deve executar mentalmente para resolver o problema sem queimar etapas;
  \item \textbf{\texttt{\#\# Formato de Sa{\'i}da}}: Descreve as estruturas, tabelas, campos obrigat{\'o}rios e modelos de documentos gerados como resultado final;
  \item \textbf{\texttt{\#\# Exemplos}}: Fornece ao menos um cen{\'a}rio aut{\^e}ntico de entrada e a correspondente sa{\'i}da esperada, servindo como refer{\^e}ncia para a calibra{\c c}{\~a}o do modelo (\emph{few-shot learning});
  \item \textbf{\texttt{\#\# Limita{\c c}{\~o}es}}: Circunscreve explicitamente o escopo da skill, delimitando o que ela n{\~a}o realiza e quando o professor deve buscar outros instrumentos ou an{\'a}lise humana;
  \item \textbf{\texttt{\#\# Depend{\^e}ncias}}: Lista formal de outras skills do reposit{\'o}rio (identificadas exclusivamente pelos seus slugs can{\^o}nicos) que complementam a a{\c c}{\~a}o da skill atual;
  \item \textbf{\texttt{\#\# Refer{\^e}ncias}}: Bibliografia cient{\'i}fica e normativa que ancora as t{\'e}cnicas apresentadas, formatada integralmente no padr{\~a}o ABNT.
\end{enumerate}

\section{Como o OpenCode Opera as Skills: A Ferramenta \texttt{skill}}

O \textbf{OpenCode} implementa uma arquitetura baseada em agentes com chamada de ferramentas (\emph{tool calling}). Ao ser inicializado, o execut{\'a}vel do OpenCode varre os diret{\'o}rios de extens{\~o}es e monta no \emph{prompt de sistema} uma lista sin{\'o}ptica com o identificador e a descri{\c c}{\~a}o de cada skill registrada.

\begin{figure}[htbp]
  \centering
  \begin{tcolorbox}[colback=boxbg, colframe=primary, arc=2mm, title={\bfseries Ciclo de Vida e Ativa{\c c}{\~a}o de uma Skill no OpenCode}]
    \small
    \begin{enumerate}
      \item \textbf{Declara{\c c}{\~a}o Inicial}: O OpenCode exp{\~o}e a ferramenta interna \texttt{skill(name="...")} e lista os metadados de todas as skills instaladas em \texttt{\~{}/.config/opencode/skills/};
      \item \textbf{Detec{\c c}{\~a}o de Inten{\c c}{\~a}o}: O usu{\'a}rio formula uma demanda docente (ex.: \emph{``Gere uma rubrica anal{\'i}tica para um semin{\'a}rio de f{\'i}sica experimental com 4 crit{\'e}rios e n{\'i}vel AIAS 2''});
      \item \textbf{Invoca{\c c}{\~a}o da Ferramenta}: O modelo de IA reconhece a correspond{\^e}ncia com a descri{\c c}{\~a}o da skill e executa a chamada de fun{\c c}{\~a}o:
      \begin{center}
        \texttt{[tool\_call: skill(name="ia-educacao-rubrica")]}
      \end{center}
      \item \textbf{Inje{\c c}{\~a}o Din{\^a}mica}: O OpenCode l{\^e} o conte{\'u}do completo de \texttt{skills/ia-educacao-rubrica/SKILL.md} e o anexa ao hist{\'o}rico da conversa como resultado da ferramenta;
      \item \textbf{Racioc{\'i}nio Guiado}: O modelo passa a operar sob as instru{\c c}{\~o}es do \texttt{Workflow}, respeitando rigorosamente os \texttt{Princ{\'i}pios}, o \texttt{Formato de Sa{\'i}da} e as regras pedag{\'o}gicas estipuladas.
    \end{enumerate}
  \end{tcolorbox}
  \caption{Mecanismo de ativa{\c c}{\~a}o da ferramenta \texttt{skill} no OpenCode}
  \label{fig:fluxo-skill}
\end{figure}

\begin{dicaopencode}
\textbf{Prontid{\~a}o de Execu{\c c}{\~a}o no Reposit{\'o}rio}: Ao contr{\'a}rio de ambientes que exigem o comando manual de adi{\c c}{\~a}o de pacotes externos, no OpenCode as skills disponibilizadas no diret{\'o}rio de configura{\c c}{\~a}o j{\'a} ficam imediatamente acess{\'i}veis. N{\~a}o {\'e} necess{\'a}rio executar comandos como \texttt{npx skills add}; a detec{\c c}{\~a}o e o acionamento ocorrem de maneira fluida e nativa.
\end{dicaopencode}

\section{O Grafo de Depend{\^e}ncias e a Integridade do Reposit{\'o}rio}

As 62 skills do reposit{\'o}rio n{\~a}o operam isoladamente; constituem um \textbf{grafo direcionado e ac{\'i}clico de conhecimento pedag{\'o}gico}. 

Por exemplo, a skill de constru{\c c}{\~a}o de avalia{\c c}{\~o}es (\texttt{ia-educacao-avaliacao}) aponta depend{\^e}ncias formais para \texttt{ia-educacao-bloom} (para alinhamento cognitivo), \texttt{ia-educacao-rubrica} (para crit{\'e}rios de corre{\c c}{\~a}o), \texttt{aias-consultant} (para governan{\c c}a do uso de IA) e \texttt{ia-educacao-planejamento-reverso} (para enquadramento curricular).

Para assegurar que o grafo n{\~a}o se degenere com o crescimento do acervo, o reposit{\'o}rio conta com uma su{\'i}te cont{\'i}nua de auditoria corporificada no script \texttt{audit.sh}:
\begin{itemize}
  \item \textbf{Zero Skills {\'O}rf{\~a}s}: Toda skill do acervo deve ser referenciada como depend{\^e}ncia por ao menos uma outra skill ou por um subagente, garantindo relev{\^a}ncia sist{\^e}mica;
  \item \textbf{Zero Depend{\^e}ncias Quebradas}: Qualquer refer{\^e}ncia cruzada em \texttt{\#\# Depend{\^e}ncias} deve corresponder a uma pasta real e v{\'a}lida dentro de \texttt{skills/};
  \item \textbf{Isolamento de Depend{\^e}ncias Externas}: As skills deste reposit{\'o}rio s{\~a}o completamente autocontidas no ecossistema brasileiro. Nenhuma skill referencia pacotes externos n{\~a}o controlados;
  \item \textbf{Conformidade Terminol{\'o}gica}: O script verifica a consist{\^e}ncia exata dos nomes dos 5 n{\'i}veis da Escala AIAS e a grafia institucional em refer{\^e}ncias bibliogr{\'a}ficas normativas do MEC.
\end{itemize}

\chapter{Instala{\c c}{\~a}o e Configura{\c c}{\~a}o do OpenCode.ai}
\label{cap:instalacao}

\section{Vis{\~a}o Geral do OpenCode}

O \textbf{OpenCode} (\url{https://opencode.ai}) {\'e} um agente aut{\^o}nomo de terminal e assistente de desenvolvimento de c{\'o}digo aberto concebido para operar diretamente sobre o sistema de arquivos local. Ao contr{\'a}rio de plataformas propriet{\'a}rias fechadas, o OpenCode oferece:

\begin{itemize}
  \item Total flexibilidade na escolha de provedores de modelos de linguagem (OpenAI, Anthropic, Google, Groq, DeepSeek ou modelos de c{\'o}digo aberto rodando localmente via Ollama);
  \item Interface interativa de terminal rica em recursos visuais (\emph{Terminal User Interface} -- TUI), al{\'e}m de suporte a execu{\c c}{\~a}o em lote (\emph{headless});
  \item Suporte nativo a ferramentas de inspe{\c c}{\~a}o de arquivos (\texttt{read}, \texttt{glob}, \texttt{grep}), edi{\c c}{\~a}o segura (\texttt{edit}, \texttt{write}), execu{\c c}{\~a}o de comandos de terminal (\texttt{bash}), extens{\~o}es via \emph{Model Context Protocol} (MCP) e gerenciamento de \emph{skills} modulares.
\end{itemize}

\section{Pr{\'e}-requisitos de Sistema}

O OpenCode {\'e} compat{\'i}vel com os principais sistemas operacionais modernos:
\begin{itemize}
  \item \textbf{Linux}: Qualquer distribui{\c c}{\~a}o contempor{\^a}nea com arquitetura x86\_64 ou ARM64 (Ubuntu 20.04+, Debian 11+, Fedora, Arch Linux, etc.);
  \item \textbf{macOS}: Vers{\~a}o 12 (Monterey) ou superior, tanto em processadores Apple Silicon (M1/M2/M3/M4) quanto Intel;
  \item \textbf{Windows}: Recomenda-se enfaticamente a utiliza{\c c}{\~a}o por meio do \textbf{WSL2} (\emph{Windows Subsystem for Linux}), que oferece ambiente Linux nativo com m{\'a}ximo desempenho e compatibilidade. Tamb{\'e}m {\'e} poss{\'i}vel a instala{\c c}{\~a}o direta no PowerShell via gerenciadores de pacotes Windows.
\end{itemize}

Ferramentas utilit{\'a}rias b{\'a}sicas no terminal: \texttt{curl}, \texttt{git} e uma conex{\~a}o ativa com a internet para baixa dos bin{\'a}rios e comunica{\c c}{\~a}o com as APIs.

\section{Passo a Passo de Instala{\c c}{\~a}o}

O OpenCode disponibiliza m{\'u}ltiplos m{\'e}todos de instala{\c c}{\~a}o oficiais. O usu{\'a}rio pode optar pelo que melhor se adeque ao seu ambiente de trabalho.

\subsection{M{\'e}todo 1: Script Autom{\'a}tico via cURL (Recomendado para Linux e macOS)}

O m{\'e}todo mais r{\'a}pido e direto consiste em executar o script oficial de instala{\c c}{\~a}o no terminal:

\begin{lstlisting}[language=bash, caption={Instala{\c c}{\~a}o r{\'a}pida via cURL}]
curl -fsSL https://opencode.ai/install | bash
\end{lstlisting}

O script detecta a arquitetura do processador, baixa o bin{\'a}rio est{\'a}tico otimizado, instala o execut{\'a}vel no diret{\'o}rio \texttt{\~{}/.opencode/bin/opencode} e adiciona essa localiza{\c c}{\~a}o ao arquivo de configura{\c c}{\~a}o do \emph{shell} (\texttt{\~{}/.bashrc} ou \texttt{\~{}/.zshrc}).

Ap{\'o}s a conclus{\~a}o, recarregue as vari{\'a}veis de ambiente do terminal:
\begin{lstlisting}[language=bash]
source ~/.bashrc   # ou source ~/.zshrc se estiver no zsh
\end{lstlisting}

\subsection{M{\'e}todo 2: Gerenciadores de Pacotes JavaScript (Node.js / npm)}

Caso o ambiente j{\'a} conte com o Node.js (vers{\~a}o 18 ou superior) configurado, o OpenCode pode ser instalado globalmente via \texttt{npm}, \texttt{pnpm} ou \texttt{bun}:

\begin{lstlisting}[language=bash, caption={Instala{\c c}{\~a}o global via npm, pnpm ou bun}]
# Via npm
npm install -g opencode-ai

# Ou via pnpm
pnpm add -g opencode-ai

# Ou via bun
bun add -g opencode-ai
\end{lstlisting}

\subsection{M{\'e}todo 3: Homebrew (macOS e Linux)}

Usu{\'a}rios do gerenciador \textbf{Homebrew} podem instalar com um {\'u}nico comando:

\begin{lstlisting}[language=bash, caption={Instala{\c c}{\~a}o via Homebrew}]
brew install opencode
\end{lstlisting}

\subsection{M{\'e}todo 4: Windows Nativo (Scoop ou Chocolatey)}

Para instala{\c c}{\~a}o direta no Windows sem WSL, abra o PowerShell como Administrador:

\begin{lstlisting}[language=bash, caption={Instala{\c c}{\~a}o no Windows via Scoop ou Chocolatey}]
# Via Scoop
scoop install opencode

# Ou via Chocolatey
choco install opencode
\end{lstlisting}

\subsection{Verifica{\c c}{\~a}o da Instala{\c c}{\~a}o}

Para confirmar se o execut{\'a}vel est{\'a} operando corretamente, consulte a vers{\~a}o no terminal:

\begin{lstlisting}[language=bash]
opencode --version
\end{lstlisting}
A sa{\'i}da dever{\'a} exibir o n{\'u}mero da vers{\~a}o instalada (por exemplo, \texttt{1.18.29}).

\section{Configura{\c c}{\~a}o de Provedores e Modelos de IA}

O OpenCode gerencia credenciais de forma segura por meio do comando \texttt{providers}.

\begin{lstlisting}[language=bash, caption={Gerenciamento de provedores e credenciais}]
# Listar os provedores configurados atualmente
opencode providers list

# Iniciar o assistente interativo de login em um provedor
opencode providers login
\end{lstlisting}

Durante a autentica{\c c}{\~a}o, o terminal solicitar{\'a} a chave de API (\emph{API Key}) do servi{\c c}o escolhido:
\begin{itemize}
  \item \textbf{Anthropic}: Chave da API Claude (\texttt{sk-ant-...}), recomendada para modelos com forte racioc{\'i}nio instrucional e reflexivo (e.g., Claude 3.5 Sonnet ou Claude 3.7 Sonnet);
  \item \textbf{OpenAI}: Chave da API OpenAI (\texttt{sk-...}) para a fam{\'i}lia GPT-4o e modelos da s{\'e}rie \emph{o1};
  \item \textbf{Google Gemini}: Chave do Google AI Studio para acesso ao Gemini 2.0 Flash e Gemini 1.5 Pro;
  \item \textbf{Modelos Locais (Ollama)}: Para institui{\c c}{\~o}es com pol{\'i}ticas restritas de privacidade de dados que n{\~a}o admitem tr{\'a}fego de dados discentes para nuvens comerciais, {\'e} poss{\'i}vel conectar o OpenCode a inst{\^a}ncias locais do Ollama (\url{http://localhost:11434}).
\end{itemize}

\section{Instala{\c c}{\~a}o e Sincroniza{\c c}{\~a}o das Skills no OpenCode}

Para que as 62 skills do reposit{\'o}rio fiquem dispon{\'i}veis no seu ambiente local do OpenCode, as pastas de skills precisam residir no diret{\'o}rio de extens{\~o}es do usu{\'a}rio:
\begin{center}
  \texttt{\~{}/.config/opencode/skills/<slug>/SKILL.md}
\end{center}

O reposit{\'o}rio prov{\^e} um script de instala{\c c}{\~a}o automatizado que realiza a c{\'o}pia e organiza{\c c}{\~a}o de forma segura.

\subsection{Executando a Instala{\c c}{\~a}o}

Na raiz do reposit{\'o}rio \texttt{skills-ia-educacao}, execute:

\begin{lstlisting}[language=bash, caption={Instala{\c c}{\~a}o automatizada das skills no OpenCode}]
# 1. Simular a instalacao para inspecionar os caminhos de destino (dry-run)
./install-skills.sh --opencode --dry-run

# 2. Efetivar a instalacao de todas as 62 skills
./install-skills.sh --opencode
\end{lstlisting}

\subsection{Confer{\^e}ncia da Instala{\c c}{\~a}o}

Para certificar-se de que todas as skills foram copiadas com sucesso:

\begin{lstlisting}[language=bash, caption={Verificando as skills instaladas}]
ls ~/.config/opencode/skills/ | grep ia-educacao | wc -l
\end{lstlisting}
O retorno dever{\'a} apontar 61 skills com o prefixo institucional (mais a skill especial \texttt{aias-consultant}, totalizando 62 unidades).

\section{Primeiros Comandos no OpenCode}

\subsection{Iniciando a Interface Interativa (TUI)}

Basta digitar o comando no terminal dentro do diret{\'o}rio do seu projeto acad{\^e}mico ou disciplina:

\begin{lstlisting}[language=bash]
opencode
\end{lstlisting}

Uma interface limpa e intuitiva ser{\'a} renderizada no terminal, permitindo enviar mensagens, aprovar a{\c c}{\~o}es de leitura/escrita e selecionar o modelo desejado em tempo real.

\subsection{Execu{\c c}{\~a}o Direta de Comandos (\emph{Headless Mode})}

Para instru{\c c}{\~o}es r{\'a}pidas ou integra{\c c}{\~a}o com scripts do professor, utilize o subcomando \texttt{run}:

\begin{lstlisting}[language=bash, caption={Exemplo de execu{\c c}{\~a}o direta com skill}]
opencode run "Utilize a skill ia-educacao-rubrica para estruturar uma rubrica analitica para uma apresentacao oral de estagio com nivel AIAS 1."
\end{lstlisting}

\subsection{Atualiza{\c c}{\~a}o do OpenCode}

Para manter o OpenCode sempre na {\'u}ltima vers{\~a}o est{\'a}vel:
\begin{lstlisting}[language=bash]
opencode upgrade
\end{lstlisting}

\chapter{Fluxo de Trabalho na Interface Gr{\'a}fica e Desktop}
\label{cap:desktop}

\begin{objetivos}
\begin{itemize}
  \item Iniciar a interface gr{\'a}fica do OpenCode no diret{\'o}rio da sua disciplina;
  \item Reconhecer as cinco regi{\~o}es da interface (sess{\~o}es, arquivos, di{\'a}logo, aprova{\c c}{\~a}o de a{\c c}{\~o}es e seletor de modelos);
  \item Executar o ciclo completo de trabalho docente na interface: abrir, selecionar, inspecionar, aprovar e exportar;
  \item Avaliar, opcionalmente, o uso do controle de vers{\~a}o com Git para proteger seus materiais.
\end{itemize}
\end{objetivos}

\section{Para al{\'e}m do Terminal: A Interface Visual do OpenCode}

Embora o OpenCode tenha nascido com foco em ambientes de linha de comando (\emph{Command Line Interface} -- CLI), o ecossistema evoluiu para oferecer suporte pleno a ambientes gr{\'a}ficos interativos. Essa caracter{\'i}stica {\'e} especialmente valiosa para o corpo docente, uma vez que professores de diferentes {\'a}reas do conhecimento nem sempre possuem familiaridade com a manipula{\c c}{\~a}o direta de comandos no \emph{shell}.

O OpenCode disponibiliza dois modos gr{\'a}ficos complementares:
\begin{itemize}
  \item \textbf{OpenCode Web UI}: Servidor local embutido iniciado por comando de terminal que disponibiliza uma interface completa no navegador web (\emph{browser});
  \item \textbf{OpenCode Desktop App}: Aplicativo nativo empacotado que integra o motor do agente a uma janela aut{\^o}noma com gerenciamento de janelas e notifica{\c c}{\~o}es do sistema.
\end{itemize}

\section{Iniciando a Interface Web e Desktop}

Para inicializar a interface gr{\'a}fica no diret{\'o}rio da sua disciplina ou material did{\'a}tico, basta executar no terminal:

\begin{lstlisting}[language=bash, caption={Iniciando a interface visual do OpenCode}]
# Inicia o servidor local e abre automaticamente o navegador padrao
opencode web

# Ou apontando diretamente para a pasta de uma disciplina especifica
opencode web ~/documentos/disciplinas/calculo-1/
\end{lstlisting}

O terminal exibir{\'a} o endere{\c c}o de acesso local (por padr{\~a}o, \url{http://127.0.0.1:4096} ou porta din{\^a}mica equivalente), abrindo instantaneamente uma sess{\~a}o gr{\'a}fica interativa.

\begin{dicaopencode}
\textbf{Acesso Remoto ou em Rede Local}: {\'E} poss{\'i}vel iniciar o servidor definindo portas e dom{\'i}nios locais customizados (por exemplo, \texttt{opencode web -{}-port 8080 -{}-hostname 0.0.0.0}) para que o professor acesse seu assistente a partir de um tablet ou outro computador conectado {\`a} rede da universidade.
\end{dicaopencode}

\section{Anatomia da Interface Gr{\'a}fica}

A interface gr{\'a}fica do OpenCode foi concebida para priorizar a clareza, a ergonomia e o controle humano sobre as a{\c c}{\~o}es aut{\^o}nomas do agente. A tela organiza-se em cinco regi{\~o}es fundamentais:

\begin{enumerate}
  \item \textbf{Barra Lateral de Hist{\'o}rico e Sess{\~o}es (\emph{Sidebar})}: Permite criar, alternar e renomear m{\'u}ltiplas sess{\~o}es de trabalho. O professor pode manter uma sess{\~a}o dedicada ao ``Planejamento da Aula 04'', outra para ``Elabora{\c c}{\~a}o da Prova P1'' e uma terceira para ``Revis{\~a}o de Artigo'';
  \item \textbf{Explorador de Arquivos e Contexto}: Exibe a {\'a}rvore de diret{\'o}rios do projeto. Permite arrastar arquivos (como PDFs normativos, ementas em \texttt{.txt} ou tabelas de notas) diretamente para a {\'a}rea de conversa;
  \item \textbf{Painel Central de Di{\'a}logo e Renderiza{\c c}{\~a}o}: Apresenta as mensagens com formata{\c c}{\~a}o rica em Markdown, suporte a f{\'o}rmulas matem{\'a}ticas em \LaTeX, tabelas din{\^a}micas e realce colorido de sintaxe de c{\'o}digo;
  \item \textbf{Painel de Auditoria e Aprova{\c c}{\~a}o de A{\c c}{\~o}es (\emph{Diff Viewer})}: Quando um subagente prop{\~o}e criar ou editar um documento no disco, a interface exibe uma visualiza{\c c}{\~a}o comparativa em duas colunas (\emph{diff} com cores verde para inclus{\~o}es e vermelho para remo{\c c}{\~o}es). Nenhuma modifica{\c c}{\~a}o {\'e} persistida sem a autoriza{\c c}{\~a}o expl{\'i}cita do usu{\'a}rio no bot{\~a}o de confirma{\c c}{\~a}o;
  \item \textbf{Seletor Superior de Modelos e Agentes}: Um menu suspenso permite alternar rapidamente entre provedores (Anthropic, OpenAI, modelos gratuitos do OpenCode ou Ollama local) e definir se a tarefa ser{\'a} conduzida pelo agente geral ou por um subagente especializado (e.g., \texttt{@construtor-de-avaliacoes}). O detalhamento completo dos seis subagentes do acervo encontra-se no Cap{\'i}tulo~\ref{cap:agentes-skills}.
\end{enumerate}

\section{O Fluxo de Trabalho Docente na Interface Gr{\'a}fica}

Para ilustrar a din{\^a}mica de trabalho na interface desktop, a Figura~\ref{fig:fluxo-desktop} resume o ciclo iterativo t{\'i}pico percorrido pelo professor ao preparar materiais pedag{\'o}gicos.

\begin{figure}[htbp]
  \centering
  \begin{tcolorbox}[colback=boxbg, colframe=primary, arc=2mm, title={\bfseries Ciclo de Trabalho Docente no OpenCode Desktop}]
    \small
    \begin{enumerate}
      \item \textbf{Abertura do Espa{\c c}o de Trabalho}: Inicia-se o OpenCode no diret{\'o}rio da disciplina contendo as ementas e objetivos institucionais;
      \item \textbf{Sele{\c c}{\~a}o do Modelo e do Subagente}: No menu superior, o docente seleciona o agente desejado (ex.: \texttt{@planejador-pedagogico}) e o modelo apropriado (gratuito ou avan{\c c}ado);
      \item \textbf{Envio de Insumos}: O docente formula a necessidade educacional e anexa o arquivo da ementa com um clique;
      \item \textbf{Inspe{\c c}{\~a}o do Racioc{\'i}nio e Chamada de Skills}: A interface destaca visualmente a ativa{\c c}{\~a}o das skills correspondentes (ex.: \texttt{/ia-educacao-bloom}, \texttt{/aias-consultant});
      \item \textbf{Aprova{\c c}{\~a}o Visual de Arquivos}: Ap{\'o}s revisar a pr{\'e}-visualiza{\c c}{\~a}o do plano ou rubrica gerada, o docente clica em \emph{Aprovar e Salvar};
      \item \textbf{Exporta{\c c}{\~a}o Final}: O artefato {\'e} salvo como arquivo \texttt{.md}, \texttt{.tex} ou \texttt{.html}, pronto para upload no ambiente virtual de aprendizagem (Moodle/SIGAA) ou impress{\~a}o.
    \end{enumerate}
  \end{tcolorbox}
  \caption{Fluxo de trabalho integrado na interface gr{\'a}fica do OpenCode}
  \label{fig:fluxo-desktop}
\end{figure}

\begin{atencaopedagogica}
\textbf{Seguran{\c c}a e Controle de Vers{\~a}o (opcional)}: O uso do OpenCode n{\~a}o exige conhecimento pr{\'e}vio de Git, e a interface gr{\'a}fica j{\'a} oferece uma primeira camada de prote{\c c}{\~a}o: nenhum arquivo {\'e} alterado sem a sua aprova{\c c}{\~a}o visual no painel de \emph{diff}. Para docentes que desejam uma camada adicional de seguran{\c c}a -- como recuperar vers{\~o}es anteriores de um plano de aula ou de uma prova ap{\'o}s uma edi{\c c}{\~a}o equivocada --, {\'e} poss{\'i}vel manter o diret{\'o}rio do curso sob controle de vers{\~a}o Git. Trata-se de uma medida recomendada, por{\'e}m facultativa, que pode ser implementada progressivamente. O passo a passo completo, pensado para quem nunca utilizou Git, encontra-se no Cap{\'i}tulo~\ref{cap:instalacao}, se{\c c}{\~a}o ``Controle de Vers{\~a}o com Git para Iniciantes''.
\end{atencaopedagogica}

\chapter{Escolha de Modelos Gratuitos no OpenCode}
\label{cap:modelos-gratuitos}

\begin{objetivos}
\begin{itemize}
  \item Conhecer os sete modelos gratuitos dispon{\'i}veis no OpenCode e os perfis de cada um;
  \item Selecionar o modelo mais adequado a cada tarefa pedag{\'o}gica com apoio da matriz de decis{\~a}o;
  \item Alternar entre modelos na linha de comando, na interface TUI e na interface gr{\'a}fica.
\end{itemize}
\end{objetivos}

\section{A Democratiza{\c c}{\~a}o do Acesso {\`a} Intelig{\^e}ncia Artificial}

Um dos maiores obst{\'a}culos enfrentados por professores de institui{\c c}{\~o}es p{\'u}blicas de ensino reside no custo financeiro associado {\`a}s assinaturas corporativas de IA e {\`a} cobran{\c c}a por tokens em APIs comerciais pagas em d{\'o}lar.

Para superar essa barreira de acesso e permitir que qualquer docente possa planejar aulas e construir instrumentos avaliativos sem {\^o}nus financeiro, o ecossistema \textbf{OpenCode} disponibiliza uma seleta camada de \textbf{modelos gratuitos da comunidade} (\emph{Free Tier / Community Models}). Esses modelos s{\~a}o mantidos e hospedados por grandes centros de pesquisa, cons{\'o}rcios abertos e pela infraestrutura do pr{\'o}prio ecossistema, oferecendo cotas generosas de requisi{\c c}{\~o}es di{\'a}rias.

Vale registrar que os modelos n{\~a}o s{\~a}o criados pela Anomaly (empresa desenvolvedora do OpenCode): cada um {\'e} propriedade de seu respectivo laborat{\'o}rio de origem -- por exemplo, a fam{\'i}lia \textbf{Nemotron} {\'e} desenvolvida pela \textbf{NVIDIA}, o \textbf{Ling} pelo cons{\'o}rcio aberto \textbf{inclusionAI}, o \textbf{MiMo} pela \textbf{Xiaomi} e o \textbf{Muse Spark} pela \textbf{Meta}. O OpenCode atua apenas como cliente que orquestra o acesso a esses modelos por meio dos identificadores \texttt{opencode/<modelo>}.

Neste cap{\'i}tulo, analisamos detalhadamente os sete modelos gratuitos mais destacados no OpenCode, suas propriedades arquiteturais e as tarefas pedag{\'o}gicas para as quais cada um {\'e} mais indicado. Os subagentes mencionados nas recomenda{\c c}{\~o}es (ex.: \path{@planejador-pedagogico}, \path{@construtor-de-avaliacoes}) s{\~a}o apresentados integralmente no Cap{\'i}tulo~\ref{cap:agentes-skills}.

\section{An{\'a}lise dos 7 Modelos Gratuitos do OpenCode}

\subsection{1. Big Pickle}

\begin{itemize}
  \item \textbf{Identificador}: \path{opencode/big-pickle};
  \item \textbf{Perfil Arquitetural}: Modelo concebido com grande janela de contexto e forte capacidade de s{\'i}ntese enciclop{\'e}dica;
  \item \textbf{For{\c c}as Pedag{\'o}gicas}: Destaca-se na leitura e sumariza{\c c}{\~a}o de longos documentos normativos, projetos pedag{\'o}gicos de curso (PPCs), planos de desenvolvimento institucional (PDIs) e ementas extensas;
  \item \textbf{Quando Usar}: Ideal para a fase inicial de diagn{\'o}stico com o subagente \path{@planejador-pedagogico}, quando o professor precisa cruzar a ementa oficial da disciplina com as compet{\^e}ncias do egresso estabelecidas pelas DCNs (Diretrizes Curriculares Nacionais).
\end{itemize}

\subsection{2. Ling 3.0 Flash Fin Free}

\begin{itemize}
  \item \textbf{Identificador}: \path{opencode/ling-3.0-flash-fin-free};
  \item \textbf{Perfil Arquitetural}: Modelo otimizado para velocidade de infer{\^e}ncia (\emph{flash}) e estrita precis{\~a}o num{\'e}rica, l{\'o}gica e quantitativa;
  \item \textbf{For{\c c}as Pedag{\'o}gicas}: Baix{\'i}ssimo {\'i}ndice de erros aritm{\'e}ticos; excelente manipula{\c c}{\~a}o de tabelas, pontua{\c c}{\~o}es ponderadas e f{\'o}rmulas de avalia{\c c}{\~a}o;
  \item \textbf{Quando Usar}: Perfeito para o c{\'a}lculo e calibra{\c c}{\~a}o de notas em rubricas anal{\'i}ticas complexas, montagem de blueprints de prova com distribui{\c c}{\~a}o de pontos e problemas de engenharia econ{\^o}mica, estat{\'i}stica ou matem{\'a}tica financeira.
\end{itemize}

\subsection{3. Mimo V2.5 Free}

\begin{itemize}
  \item \textbf{Identificador}: \path{opencode/mimo-v2.5-free};
  \item \textbf{Perfil Arquitetural}: Modelo vers{\'a}til, leve e equilibrado, com excelente alinhamento {\`a}s instru{\c c}{\~o}es fornecidas em l{\'i}ngua portuguesa;
  \item \textbf{For{\c c}as Pedag{\'o}gicas}: Reda{\c c}{\~a}o fluida, tom equilibrado (nem rob{\'o}tico, nem informal), facilidade para organizar respostas em t{\'o}picos e alta fidelidade aos esquemas de frontmatter das skills;
  \item \textbf{Quando Usar}: O modelo ``faz-tudo'' para a rotina di{\'a}ria. Recomendado para a reda{\c c}{\~a}o de enunciados de quest{\~o}es discursivas, esbo{\c c}o de cronogramas semanais e media{\c c}{\~a}o de atividades ativas em sala de aula.
\end{itemize}

\subsection{4. Muse Spark 1.2 Free}

\begin{itemize}
  \item \textbf{Identificador}: \path{opencode/muse-spark-1.2-contributor-free};
  \item \textbf{Perfil Arquitetural}: Modelo com calibra{\c c}{\~a}o voltada para idea{\c c}{\~a}o did{\'a}tica, est{\'i}mulo criativo e gera{\c c}{\~a}o de hip{\'o}teses;
  \item \textbf{For{\c c}as Pedag{\'o}gicas}: Capacidade de sugerir analogias pedag{\'o}gicas v{\'a}lidas, met{\'a}foras explicativas e cen{\'a}rios originais do mundo real para problematiza{\c c}{\~a}o contextualizada;
  \item \textbf{Quando Usar}: Na fase de \emph{brainstorming} com a skill \path{ia-educacao-design-problema}, formula{\c c}{\~a}o de situa{\c c}{\~o}es-gatilho para PBL (\emph{Problem-Based Learning}) e din{\^a}micas de debate acad{\^e}mico.
\end{itemize}

\subsection{5. Muse Spark 1.3 Free}

\begin{itemize}
  \item \textbf{Identificador}: \path{opencode/muse-spark-1.3-contributor-free};
  \item \textbf{Perfil Arquitetural}: Evolu{\c c}{\~a}o direta da fam{\'i}lia Muse Spark, combinando a criatividade did{\'a}tica da vers{\~a}o anterior com rigor aprimorado na estrutura de sa{\'i}da;
  \item \textbf{For{\c c}as Pedag{\'o}gicas}: Gera{\c c}{\~a}o de c{\'o}digo \LaTeX{} e tabelas Markdown perfeitamente alinhadas, com consist{\^e}ncia sem{\^a}ntica e rara ocorr{\^e}ncia de alucina{\c c}{\~o}es sint{\'a}ticas;
  \item \textbf{Quando Usar}: Constru{\c c}{\~a}o de rubricas anal{\'i}ticas tabuladas (\path{ia-educacao-rubrica}) e confec{\c c}{\~a}o de exames formais que exigem diagramas ou equa{\c c}{\~o}es matem{\'a}ticas rigorosas.
\end{itemize}

\subsection{6. Nemotron 3 Ultra Free}

\begin{itemize}
  \item \textbf{Identificador}: \path{opencode/nemotron-3-ultra-free};
  \item \textbf{Perfil Arquitetural}: Modelo denso de alt{\'i}ssima capacidade baseado na arquitetura aberta da NVIDIA, com bilh{\~o}es de par{\^a}metros otimizados para racioc{\'i}nio l{\'o}gico avan{\c c}ado;
  \item \textbf{For{\c c}as Pedag{\'o}gicas}: Compreens{\~a}o profunda de c{\'o}digo-fonte (C, C++, Python, Java, Rust), arquitetura de software, teoremas matem{\'a}ticos e demonstra{\c c}{\~o}es formais;
  \item \textbf{Quando Usar}: Avalia{\c c}{\~a}o de projetos t{\'e}cnicos em Ci{\^e}ncia da Computa{\c c}{\~a}o e Engenharias, resolu{\c c}{\~a}o anal{\'i}tica de exerc{\'i}cios de f{\'i}sica avan{\c c}ada e auditoria cr{\'i}tica de gabaritos com o subagente \path{@construtor-de-avaliacoes}.
\end{itemize}

\subsection{7. Nemotron 3.5 Lightning Free}

\begin{itemize}
  \item \textbf{Identificador}: \path{opencode/nemotron-3.5-lightning-free};
  \item \textbf{Perfil Arquitetural}: Variante de resposta quase instant{\^a}nea (\emph{lightning}), balanceando a for{\c c}a l{\'o}gica da arquitetura Nemotron com lat{\^e}ncia ultrabaixa;
  \item \textbf{For{\c c}as Pedag{\'o}gicas}: Devolutivas r{\'a}pidas em fra{\c c}{\~o}es de segundo, mantendo coer{\^e}ncia em disciplinas STEM;
  \item \textbf{Quando Usar}: Sess{\~o}es interativas em tempo real (como a gera{\c c}{\~a}o dinâmica de testes conceituais durante a aula para \emph{Peer Instruction} via \path{ia-educacao-peer-instruction}) e devolu{\c c}{\~a}o r{\'a}pida de feedback para m{\'u}ltiplos estudantes.
\end{itemize}

\begin{itemize}
  \item \textbf{Perfil Arquitetural}: Variante de resposta quase instant{\^a}nea (\emph{lightning}), balanceando a for{\c c}a l{\'o}gica da arquitetura Nemotron com lat{\^e}ncia ultrabaixa;
  \item \textbf{For{\c c}as Pedag{\'o}gicas}: Devolutivas r{\'a}pidas em fra{\c c}{\~o}es de segundo, mantendo coer{\^e}ncia em disciplinas STEM;
  \item \textbf{Quando Usar}: Sess{\~o}es interativas em tempo real (como a gera{\c c}{\~a}o dinâmica de testes conceituais durante a aula para \emph{Peer Instruction} via \path{ia-educacao-peer-instruction}) e devolu{\c c}{\~a}o r{\'a}pida de feedback para m{\'u}ltiplos estudantes.
\end{itemize}

\section{Matriz de Decis{\~a}o Docente: Qual Modelo Escolher?}

A Tabela~\ref{tab:modelos-matriz} resume a recomenda{\c c}{\~a}o de uso de cada modelo gratuito conforme a necessidade pedag{\'o}gica imediata do professor.

\begin{table}[htbp]
  \centering
  \small
  \caption{Matriz de Sele{\c c}{\~a}o de Modelos Gratuitos do OpenCode}
  \label{tab:modelos-matriz}
  \begin{tabularx}{\linewidth}{p{4.8cm}p{4.2cm}X}
    \toprule
    \textbf{Tarefa Pedag{\'o}gica} & \textbf{Modelo Recomendado} & \textbf{Identificador no OpenCode} \\
    \midrule
    An{\'a}lise de PPCs e Ementas Longas & Big Pickle & \path{opencode/big-pickle} \\
    \addlinespace
    C{\'a}lculo de Notas e Pesos & Ling 3.0 Flash Fin Free & \path{opencode/ling-3.0-flash-fin-free} \\
    \addlinespace
    Reda{\c c}{\~a}o de Aulas e Exerc{\'i}cios & Mimo V2.5 Free & \path{opencode/mimo-v2.5-free} \\
    \addlinespace
    Problemas e Estudos de Caso & Muse Spark 1.2 Free & \path{opencode/muse-spark-1.2-contributor-free} \\
    \addlinespace
    Rubricas e Formata{\c c}{\~a}o \LaTeX{} & Muse Spark 1.3 Free & \path{opencode/muse-spark-1.3-contributor-free} \\
    \addlinespace
    Demonstra{\c c}{\~o}es e C{\'o}digo STEM & Nemotron 3 Ultra Free & \path{opencode/nemotron-3-ultra-free} \\
    \addlinespace
    Sess{\~o}es em Sala e Feedback R{\'a}pido & Nemotron 3.5 Lightning Free & \path{opencode/nemotron-3.5-lightning-free} \\
    \bottomrule
  \end{tabularx}
\end{table}

\section{Como Alternar entre Modelos no OpenCode}

O professor pode selecionar ou alternar o modelo em uso de tr{\^e}s formas simples:

\begin{enumerate}
  \item \textbf{Ao Iniciar o OpenCode}: Passando o identificador pela flag \texttt{-m}:
  \begin{lstlisting}[language=bash]
  opencode -m opencode/nemotron-3-ultra-free
  \end{lstlisting}
  \item \textbf{Na Interface Interativa de Terminal (TUI)}: Pressionando a tecla de atalho \texttt{Ctrl+M} (ou digitando \texttt{/model}) e selecionando o modelo na lista;
  \item \textbf{Na Interface Web/Desktop}: Clicando no menu suspenso no canto superior direito da janela e escolhendo o modelo gratuito correspondente com o mouse.
\end{enumerate}

\chapter{{\'E}tica, Limita{\c c}{\~o}es e Transpar{\^e}ncia}
\label{cap:etica-transparencia}

\begin{objetivos}
\begin{itemize}
  \item Fundamentar o uso pedag{\'o}gico da IA nos princ{\'i}pios de centralidade humana, transpar{\^e}ncia, equidade e privacidade;
  \item Aplicar o princ{\'i}pio da supervis{\~a}o humana em todo o ciclo avaliativo;
  \item Observar os protocolos da LGPD no tratamento de dados discentes;
  \item Reconhecer as limita{\c c}{\~o}es intr{\'i}nsecas dos modelos e mitig{\'a}-las com a t{\'e}cnica CoVe.
\end{itemize}
\end{objetivos}

\section{Fundamenta{\c c}{\~a}o {\'E}tica na Educa{\c c}{\~a}o Superior}

A integra{\c c}{\~a}o de sistemas de intelig{\^e}ncia artificial na educa{\c c}{\~a}o n{\~a}o constitui uma quest{\~a}o estritamente t{\'e}cnica, mas primordialmente {\'e}tica, social e pol{\'i}tica. Como preconizam o \emph{Referencial do MEC} \cite{brasil2026referencial} e as orienta{\c c}{\~o}es da UNESCO \cite{unesco2021guidance}, a tecnologia deve servir {\`a} emancipa{\c c}{\~a}o humana, {\`a} equidade de oportunidades e ao florescimento do pensamento cr{\'i}tico, nunca {\`a} vigil{\^a}ncia punitiva ou {\`a} desumaniza{\c c}{\~a}o do percurso formativo.

No {\^a}mbito deste acervo, quatro princ{\'i}pios {\'e}ticos estruturantes orientam cada skill e subagente (os seis subagentes do reposit{\'o}rio s{\~a}o detalhados no Cap{\'i}tulo~\ref{cap:agentes-skills}):
\begin{enumerate}
  \item \textbf{Centralidade e Ag{\^e}ncia Humana}: A intelig{\^e}ncia artificial atua como co-planejadora e ferramenta de amplia{\c c}{\~a}o cognitiva; o discernimento pedag{\'o}gico e a autoridade final pertencem incondicionalmente ao docente;
  \item \textbf{Transpar{\^e}ncia e Explicabilidade}: Regras claras, pol{\'i}ticas de avalia{\c c}{\~a}o compreens{\'i}veis e declara{\c c}{\~a}o aberta sobre como e quando a IA {\'e} utilizada em sala de aula;
  \item \textbf{Equidade e Inclus{\~a}o}: Garantia de que a ado{\c c}{\~a}o da tecnologia n{\~a}o amplie desigualdades sociais ou prejudique estudantes com menor acesso a recursos digitais avan{\c c}ados;
  \item \textbf{Privacidade e Salvaguarda de Dados}: Tratamento rigoroso de informa{\c c}{\~o}es discentes em conson{\^a}ncia com a legisla{\c c}{\~a}o nacional de prote{\c c}{\~a}o de dados.
\end{enumerate}

\section{Supervis{\~a}o Humana Obrigat{\'o}ria (\emph{Human-in-the-Loop})}

A automa{\c c}{\~a}o cega de processos avaliativos {\'e} expressamente vedada pelas diretrizes deste reposit{\'o}rio (conforme formalizado na skill \path{ia-educacao-supervisao-humana}).

\begin{atencaopedagogica}
\textbf{Veda{\c c}{\~a}o a Decis{\~o}es Automatizadas Punitivas}: {\'E} terminantemente contr{\'a}rio {\`a} {\'e}tica acad{\^e}mica reprovar, penalizar ou atribuir conceitos finais a um estudante com base exclusiva em relat{\'o}rios automatizados gerados por IA ou por supostos ``detectores de pl{\'a}gio por IA'' (os quais apresentam elevados {\'i}ndices de falsos-positivos, especialmente contra estudantes n{\~a}o nativos ou neurodivergentes). Toda avalia{\c c}{\~a}o requer o escrut{\'i}nio reflexivo e a valida{\c c}{\~a}o do professor.
\end{atencaopedagogica}

Os subagentes avaliativos (como o \path{@aplicador-de-rubricas}) geram \emph{propostas formativas} de parecer. Cabe ao docente revisar o texto, ajustar a pontua{\c c}{\~a}o, acolher as particularidades contextuais do estudante e assinar a avalia{\c c}{\~a}o final.

\section{Privacidade Discente e Conformidade com a LGPD}

No Brasil, a Lei Geral de Prote{\c c}{\~a}o de Dados Pessoais (Lei n{\textordmasculine} 13.709/2018 -- LGPD) imp{\~o}e salvaguardas r{\'i}gidas sobre o tratamento de dados de estudantes. Ao utilizar o OpenCode em disciplinas universit{\'a}rias, o corpo docente deve seguir os seguintes protocolos:

\begin{itemize}
  \item \textbf{Anonimiza{\c c}{\~a}o Prévia Obrigat{\'o}ria}: Antes de submeter trabalhos de estudantes para an{\'a}lise ou gera{\c c}{\~a}o de feedback por modelos de nuvem, remova nomes completos, n{\'u}meros de matr{\'i}cula, CPFs, endere{\c c}os eletr{\^o}nicos e quaisquer metadados que permitam a re-identifica{\c c}{\~a}o do indiv{\'i}duo;
  \item \textbf{Vedada a Inser{\c c}{\~a}o de Dados Sens{\'i}veis}: Jamais insira em \emph{prompts} informa{\c c}{\~o}es sobre laudos m{\'e}dicos discentes, condi{\c c}{\~o}es socioecon{\^o}micas vulner{\'a}veis, cren{\c c}as religiosas ou processos disciplinares em andamento;
  \item \textbf{Ado{\c c}{\~a}o de Modelos Locais}: Para dados de pesquisas acad{\^e}micas com exig{\^e}ncia de sigilo {\'e}tico (aprovadas por Comit{\^e} de {\'E}tica em Pesquisa -- CEP), configure o OpenCode para operar com modelos executados 100\% localmente via Ollama (\url{http://localhost:11434}), eliminando qualquer tr{\'a}fego de dados para servidores externos.
\end{itemize}

\section{Limita{\c c}{\~o}es Intr{\'i}nsecas dos Modelos de Linguagem}

Os modelos de linguagem contempor{\^a}neos n{\~a}o s{\~a}o bancos de dados relacionais e n{\~a}o realizam racioc{\'i}nio l{\'o}gico dedutivo infal{\'i}vel. Eles s{\~a}o sistemas probabil{\'i}sticos complexos treinados para prever a sequ{\^e}ncia mais veross{\'i}mil de palavras (\emph{tokens}). Compreender suas limita{\c c}{\~o}es intr{\'i}nsecas {\'e} essencial para evitar armadilhas pedag{\'o}gicas:

\begin{enumerate}
  \item \textbf{Alucina{\c c}{\~a}o Factual (\emph{Hallucination})}: Os modelos s{\~a}o capazes de gerar cita{\c c}{\~o}es bibliogr{\'a}ficas fict{\'i}cias com autores inexistentes, inventar datas hist{\'o}ricas ou apresentar dedu{\c c}{\~o}es matem{\'a}ticas aparentemente elegantes com erros conceituais sutis;
  \item \textbf{Vi{\'e}s de Concord{\^a}ncia (\emph{Sycophancy})}: Os LLMs tendem a concordar com as premissas do usu{\'a}rio, mesmo quando o professor comete um equ{\'i}voco intencional no \emph{prompt} (por exemplo, afirmando que uma integral impr{\'o}pria diverge quando na verdade ela converge);
  \item \textbf{Vieses Algor{\'i}tmicos e Culturais}: Devido {\`a} predomin{\^a}ncia de textos em l{\'i}ngua inglesa e corpora gerados no Norte Global, os modelos podem reproduzir estere{\'o}tipos de g{\^e}nero, sub-representar autores latino-americanos e demonstrar desconhecimento da legisla{\c c}{\~a}o educacional brasileira;
  \item \textbf{Cegueira Temporal}: Os modelos possuem cortes fixos de treinamento (\emph{knowledge cutoff}), n{\~a}o tendo conhecimento imediato de leis, editais ou acontecimentos recentes sem que esses dados sejam injetados expressamente no contexto.
\end{enumerate}

\section{Mitiga{\c c}{\~a}o de Alucina{\c c}{\~o}es com a T{\'e}cnica CoVe}

Para combater as alucina{\c c}{\~o}es e assegurar o rigor cient{\'i}fico dos instrumentos did{\'a}ticos gerados, o reposit{\'o}rio incorpora nativamente a t{\'e}cnica \textbf{Chain-of-Verification (CoVe)}, implementada na skill \path{ia-educacao-verificacao}.

O processo decomp{\~o}e a valida{\c c}{\~a}o em quatro passos l{\'o}gicos independentes:
\begin{enumerate}
  \item \textbf{Gera{\c c}{\~a}o da Resposta Base}: O modelo produz o rascunho inicial do instrumento did{\'a}tico;
  \item \textbf{Planejamento de Perguntas de Checagem}: O agente identifica todas as afirma{\c c}{\~o}es factuais cr{\'i}ticas (f{\'o}rmulas, teoremas, artigos de lei) e formula perguntas objetivas de verifica{\c c}{\~a}o;
  \item \textbf{Execu{\c c}{\~a}o Independente}: As perguntas s{\~a}o respondidas de modo isolado, sem que o modelo seja induzido pela resposta anterior;
  \item \textbf{S{\'i}ntese e Corre{\c c}{\~a}o Final}: O documento final s{\'o} {\'e} entregue ap{\'o}s confrontar a vers{\~a}o inicial com as respostas validadas, expurgando imprecis{\~o}es.
\end{enumerate}

\chapter{Como Explicar Intelig{\^e}ncia Artificial aos Estudantes}
\label{cap:explicar-estudantes}

\begin{objetivos}
\begin{itemize}
  \item Estabelecer o contrato pedag{\'o}gico de transpar{\^e}ncia no primeiro dia de aula;
  \item Desconstruir os mitos do or{\'a}culo infal{\'i}vel e da ferramenta in{\'u}til junto aos estudantes;
  \item Traduzir os cinco n{\'i}veis da Escala AIAS para linguagem acess{\'i}vel \`a turma;
  \item Implantar o protocolo de declara{\c c}{\~a}o de uso de IA nas atividades avaliativas.
\end{itemize}
\end{objetivos}

\section{O Contrato Pedag{\'o}gico da Aula Inaugural}

O maior obst{\'a}culo enfrentado por estudantes em rela{\c c}{\~a}o {\`a} intelig{\^e}ncia artificial n{\~a}o {\'e} a complexidade t{\'e}cnica das ferramentas, mas a \textbf{ambiguidade normativa}. Quando a institui{\c c}{\~a}o ou o docente n{\~a}o explicitam com clareza o que {\'e} permitido e o que {\'e} vedado, cria-se um ambiente de inseguran{\c c}a, medo e desonestidade velada \cite{perkins2024aias}.

O primeiro dia de aula {\'e} o momento determinante para estabelecer o \textbf{contrato pedag{\'o}gico de transpar{\^e}ncia}. Em vez de iniciar o semestre amea{\c c}ando puni{\c c}{\~o}es com base em detectores ineficazes, o professor deve assumir a lideran{\c c}a do di{\'a}logo:

\begin{conceitochave}
\textbf{Princ{\'i}pio da Pactua{\c c}{\~a}o Aberta}: A intelig{\^e}ncia artificial faz parte da realidade profissional e acad{\^e}mica contempor{\^a}nea. O objetivo da universidade n{\~a}o {\'e} fingir que a IA n{\~a}o existe, mas ensinar o estudante a utiliz{\'a}-la com discernimento cr{\'i}tico, honestidade intelectual, rigor metodol{\'o}gico e responsabilidade {\'e}tica.
\end{conceitochave}

\section{Desmistificando a IA em Sala de Aula}

Para que os estudantes adotem uma postura madura diante das ferramentas generativas, {\'e} necess{\'a}rio desconstruir dois mitos opostos igualmente danosos: o mito do \emph{or{\'a}culo infal{\'i}vel} (crer cegamente em tudo o que a m{\'a}quina responde) e o mito da \emph{ferramenta in{\'u}til} (desprezar a tecnologia por preconceito).

Uma analogia pedag{\'o}gica amplamente recomendada (trabalhada na skill \path{ia-educacao-pensamento-critico}) {\'e} apresentar o modelo de linguagem como um \textbf{``estagi{\'a}rio hiperveloz, mas extremamente desatento''}:
\begin{itemize}
  \item Ele leu milh{\~o}es de livros e escreve com velocidade espantosa;
  \item Ele n{\~a}o tem compreens{\~a}o real do mundo e n{\~a}o possui senso moral;
  \item Quando n{\~a}o sabe uma resposta, ele prefere inventar uma mentira convincente a admitir que n{\~a}o sabe;
  \item Portanto, entregar o relat{\'o}rio do estagi{\'a}rio ao cliente sem revis{\~a}o minuciosa {\'e} neglig{\^e}ncia profissional grave. O valor do estudante est{\'a} justamente na sua capacidade de auditar, corrigir e direcionar o trabalho.
\end{itemize}

\subsection{Din{\^a}mica Pr{\'a}tica: A ``Ca{\c c}a ao Erro''}

Uma estrat{\'e}gia did{\'a}tica de alto impacto consiste em projetar na primeira semana de aula uma resposta gerada por IA sobre um t{\'o}pico b{\'a}sico da disciplina que contenha uma alucina{\c c}{\~a}o sutil proposital (uma data equivocada, um julgado citado com tese que n{\~a}o existe, um artigo de lei citado com dispositivo incorreto, um teorema com premissas omitidas ou um c{\'o}digo com vazamento de mem{\'o}ria). 

Desafiar a turma a identificar as falhas em pequenos grupos demonstra imediatamente aos discentes que o uso passivo da IA resulta em mediocridade, enquanto o conhecimento conceitual profundo {\'e} indispens{\'a}vel para validar qualquer sa{\'i}da algor{\'i}tmica.

\section{Traduzindo a Escala AIAS para os Estudantes}

A Escala AIAS n{\~a}o deve permanecer como um documento t{\'e}cnico restrito aos professores. Ela deve ser explicada em linguagem direta e afirmativa aos estudantes, conforme sugerido na Tabela~\ref{tab:aias-discente}.

Antes de apresentar a tabela, informe duas coisas separadamente: se o uso de IA generativa {\'e} permitido ou proibido e se seu uso {\'e} opcional ou obrigat{\'o}rio. Tamb{\'e}m indique as tecnologias assistivas e adapta{\c c}{\~o}es autorizadas, que n{\~a}o devem ser confundidas com o uso de IA generativa.

\begin{table}[htbp]
  \centering
  \small
  \caption{Como Explicar os N{\'i}veis AIAS para a Turma}
  \label{tab:aias-discente}
  \begin{tabularx}{\linewidth}{clX}
    \toprule
    \textbf{N{\'i}vel} & \textbf{Nome} & \textbf{Mensagem Pedag{\'o}gica ao Estudante} \\
    \midrule
    \textbf{1} & Sem IA & ``Nesta tarefa, quero avaliar seu racioc{\'i}nio e suas decis{\~o}es sem uso de IA generativa. Tecnologias assistivas e adapta{\c c}{\~o}es previamente autorizadas continuam permitidas.'' \\
    \addlinespace
    \textbf{2} & Planejamento Assistido por IA & ``Voc{\^e} pode usar a IA para debater ideias, montar esquemas e sugerir organiza{\c c}{\~o}es de t{\'o}picos. Contudo, o texto final e os c{\'o}digos devem ser redigidos integralmente por voc{\^e}.'' \\
    \addlinespace
    \textbf{3} & Colabora{\c c}{\~a}o com IA & ``Voc{\^e} pode usar a IA para sugerir corre{\c c}{\~o}es e refinamentos em trechos espec{\'i}ficos do seu trabalho, desde que documente explicitamente quais partes contaram com aux{\'i}lio da ferramenta.'' \\
    \addlinespace
    \textbf{4} & IA Integral & ``Nesta atividade, o enunciado exige o uso da IA ao longo do processo. Voc{\^e} ser{\'a} avaliado pela qualidade das decis{\~o}es, pela auditoria das respostas, pela documenta{\c c}{\~a}o do processo e pelo dom{\'i}nio cr{\'i}tico do resultado final.'' \\
    \addlinespace
    \textbf{5} & Explora{\c c}{\~a}o de IA & ``Vamos explorar solu{\c c}{\~o}es in{\'e}ditas e interdisciplinares. Voc{\^e} dirigir{\'a} o processo, registrar{\'a} as contribui{\c c}{\~o}es da IA e continuar{\'a} respons{\'a}vel pela valida{\c c}{\~a}o e pelo resultado final.'' \\
    \bottomrule
  \end{tabularx}
\end{table}

\section{O Protocolo de Transpar{\^e}ncia e Atribui{\c c}{\~a}o}

Para atividades nos n{\'i}veis AIAS~2, 3, 4 ou 5, o plano de ensino deve instituir uma cl{\'a}usula simples e transparente de \textbf{Declara{\c c}{\~a}o de Uso de IA}:

\begin{tcolorbox}[colback=white, colframe=boxborder, arc=1mm, title={Modelo de Declara{\c c}{\~a}o de Uso de IA a ser Anexada pelo Estudante}]
\small
\textbf{Declara{\c c}{\~a}o de Transpar{\^e}ncia Acad{\^e}mica}:\\
``Neste trabalho (N{\'i}vel AIAS \underline{\hspace{1cm}}), utilizei a ferramenta \underline{\hspace{3cm}} (vers{\~a}o \underline{\hspace{1.5cm}}) com a seguinte finalidade pedag{\'o}gica: \underline{\hspace{6cm}}.\\
Os principais \emph{prompts} utilizados foram arquivados e coloco-me {\`a} disposi{\c c}{\~a}o para defender oralmente qualquer etapa ou linha deste projeto. Todas as dedu{\c c}{\~o}es e dados foram validados por mim contra fontes bibliogr{\'a}ficas can{\^o}nicas.''
\end{tcolorbox}

Essa sistem{\'a}tica substitui o medo da puni{\c c}{\~a}o pela cultura da responsabilidade acad{\^e}mica, alinhando a conduta discente aos princ{\'i}pios de integridade cient{\'i}fica defendidos pelo Minist{\'e}rio da Educa{\c c}{\~a}o \cite{brasil2026referencial}.

\begin{tcolorbox}[colback=white, colframe=boxborder, arc=1mm, title={Texto pronto para o plano de ensino (copie e adapte)}]
\small
``Nesta disciplina, o uso de IA generativa segue a Escala AIAS declarada em cada enunciado. Salvo indica{\c c}{\~a}o em contr{\'a}rio, as atividades s{\~a}o N{\'i}vel AIAS \underline{\hspace{1cm}} (\underline{\hspace{4cm}}). Quando permitido, declare a ferramenta, a vers{\~a}o e a finalidade no anexo de transpar{\^e}ncia e guarde seus \emph{prompts}. Tecnologias assistivas e adapta{\c c}{\~o}es autorizadas pela institui{\c c}{\~a}o continuam permitidas e n{\~a}o se confundem com IA generativa. Em caso de d{\'u}vida, pergunte antes de entregar.''
\end{tcolorbox}

\chapter{Agentes Especializados e Cat{\'a}logo de Skills}
\label{cap:agentes-skills}

\section{A Arquitetura de Subagentes no OpenCode}

Enquanto uma \emph{skill} fornece o conhecimento procedimental para uma a{\c c}{\~a}o circunscrita, um \textbf{agente} (ou \emph{subagente}) atua como um especialista aut{\^o}nomo orientado a metas educacionais complexas, coordenando m{\'u}ltiplas skills e operando sob pol{\'i}ticas de permiss{\~a}o estritas de seguran{\c c}a.

No OpenCode, os agentes s{\~a}o declarados com frontmatter YAML e regras de permiss{\~a}o por ferramenta (\texttt{read}, \texttt{glob}, \texttt{grep}, \texttt{edit} e \texttt{bash}):

\begin{lstlisting}[language=bash, caption={Esquema de configura{\c c}{\~a}o de um subagente no OpenCode}]
---
description: Planejamento pedagogico para ensino superior com IA.
mode: subagent
model: any
permission:
  edit: deny
  bash: deny
  read: allow
  glob: allow
  grep: allow
---
\end{lstlisting}

\subsection{Sincronia Can{\^o}nica entre \texttt{agents/} e \texttt{.opencode/agents/}}

Para garantir que as instru{\c c}{\~o}es estejam sempre atualizadas, o reposit{\'o}rio mant{\'e}m duas pastas sincronizadas com verifica{\c c}{\~a}o autom{\'a}tica via \texttt{./audit.sh}:
\begin{itemize}
  \item \texttt{agents/} -- Diret{\'o}rio fonte can{\^o}nico do reposit{\'o}rio;
  \item \texttt{.opencode/agents/} -- Diret{\'o}rio lido automaticamente pelo OpenCode no espa{\c c}o de trabalho.
\end{itemize}

\section{Os Seis Subagentes do Reposit{\'o}rio}

\subsection{1. Planejador Pedag{\'o}gico (\texttt{planejador-pedagogico})}
\begin{itemize}
  \item \textbf{Foco}: Consultor s{\^e}nior de design instrucional e planejamento curricular para o ensino superior. Analisa ementas e projetos pedag{\'o}gicos de curso (PPC), sugerindo metodologias ativas e alinhamento {\`a} Escala AIAS;
  \item \textbf{Permiss{\~o}es}: Somente leitura (\texttt{edit: deny}, \texttt{bash: deny});
  \item \textbf{Skills Recomendadas}: \path{aias-consultant}, \path{ia-educacao-avaliacao}, \path{ia-educacao-rubrica}, \path{ia-educacao-planejamento-didatico}, \path{ia-educacao-integridade-academica}, \path{ia-educacao-etica}.
\end{itemize}

\subsection{2. Construtor de Avalia{\c c}{\~o}es (\texttt{construtor-de-avaliacoes})}
\begin{itemize}
  \item \textbf{Foco}: Cria{\c c}{\~a}o de instrumentos avaliativos prontos para aplica{\c c}{\~a}o (provas, exames, listas contextualizadas, matrizes e gabaritos comentados com declara{\c c}{\~a}o formal do n{\'i}vel AIAS no cabe{\c c}alho);
  \item \textbf{Permiss{\~o}es}: Leitura e escrita (\texttt{edit: allow}, \texttt{bash: allow});
  \item \textbf{Skills Orquestradas}: \path{ia-educacao-avaliacao}, \path{ia-educacao-rubrica}, \path{ia-educacao-design-problema}, \path{ia-educacao-bloom}, \path{ia-educacao-planejamento-reverso}.
\end{itemize}

\subsection{3. Gerador de Quest{\~o}es Bloom (\texttt{gerador-de-questoes-bloom})}
\begin{itemize}
  \item \textbf{Foco}: Formula{\c c}{\~a}o de itens alinhados {\`a} Taxonomia Revisada de Bloom. Em quest{\~o}es de m{\'u}ltipla escolha, calibra distratores baseados em equ{\'i}vocos conceituais t{\'i}picos (\emph{misconceptions}) dos estudantes;
  \item \textbf{Permiss{\~o}es}: Leitura e escrita (\texttt{edit: allow}, \texttt{bash: allow});
  \item \textbf{Skills Orquestradas}: \path{ia-educacao-bloom}, \path{ia-educacao-mcq}, \path{ia-educacao-banco-questoes}, \path{ia-educacao-verificacao}, \path{aias-consultant}.
\end{itemize}

\subsection{4. Aplicador de Rubricas (\texttt{aplicador-de-rubricas})}
\begin{itemize}
  \item \textbf{Foco}: Avalia{\c c}{\~a}o formativa de trabalhos e c{\'o}digos de estudantes contra rubricas anal{\'i}ticas pr{\'e}-estabelecidas, emitindo notas descritivas por crit{\'e}rio e parecer sobre conformidade AIAS;
  \item \textbf{Permiss{\~o}es}: Leitura e escrita (\texttt{edit: allow}, \texttt{bash: allow});
  \item \textbf{Skills Orquestradas}: \path{ia-educacao-rubrica}, \path{ia-educacao-avaliacao}, \path{ia-educacao-bloom}, \path{aias-consultant}.
\end{itemize}

\subsection{5. Coach de Feedback Formativo (\texttt{coach-de-feedback-formativo})}
\begin{itemize}
  \item \textbf{Foco}: Elabora{\c c}{\~a}o de feedback formativo tridimensional (\emph{Feed Up}, \emph{Feed Back}, \emph{Feed Forward}) com calibra{\c c}{\~a}o inclusiva para estudantes t{\'i}picos, com dislexia e com TDAH;
  \item \textbf{Permiss{\~o}es}: Leitura e escrita (\texttt{edit: allow}, \texttt{bash: allow});
  \item \textbf{Skills Orquestradas}: \path{ia-educacao-feedback}, \path{ia-educacao-verificacao}, \path{ia-educacao-dua}, \path{ia-educacao-aprendizagem-ativa}, \path{ia-educacao-rubrica}.
\end{itemize}

\subsection{6. Artes{\~a}o de Skills (\texttt{artesao-de-skills})}
\begin{itemize}
  \item \textbf{Foco}: Engenheiro de prompt e guardi{\~a}o da integridade do reposit{\'o}rio. Cria novas skills, refatora m{\'o}dulos existentes e audita o grafo com \texttt{audit.sh};
  \item \textbf{Permiss{\~o}es}: Leitura e escrita (\texttt{edit: allow}, \texttt{bash: allow}).
\end{itemize}

\begin{table}[htbp]
  \centering
  \footnotesize
  \caption{S{\'i}ntese dos Subagentes Pedag{\'o}gicos para OpenCode}
  \label{tab:agentes-sintese}
  \begin{tabularx}{\linewidth}{>{\raggedright\arraybackslash}p{4.2cm}cp{3.8cm}X}
    \toprule
    \textbf{Subagente} & \textbf{Permiss{\~a}o} & \textbf{Insumos Principais} & \textbf{Entreg{\'a}veis Chave} \\
    \midrule
    \path{planejador-pedagogico} & Somente Leitura & Ementas, PPCs, carga hor{\'a}ria e objetivos & Plano de curso, cronograma e diretrizes AIAS \\
    \addlinespace
    \path{construtor-de-avaliacoes} & Leitura/Escrita & Objetivos de aprendizagem e n{\'i}vel AIAS & Provas, exerc{\'i}cios e gabaritos comentados \\
    \addlinespace
    \path{gerador-de-questoes-bloom} & Leitura/Escrita & T{\'o}pico t{\'e}cnico, n{\'i}vel de Bloom e AIAS & Quest{\~o}es discursivas e MCQs com distratores \\
    \addlinespace
    \path{aplicador-de-rubricas} & Leitura/Escrita & Entrega discente e rubrica declarada & Notas por crit{\'e}rio, parecer e an{\'a}lise de vi{\'e}s \\
    \addlinespace
    \path{coach-de-feedback-formativo} & Leitura/Escrita & Produ{\c c}{\~a}o discente e perfil neurodivergente & Feedback formativo adaptado (\emph{Feed Up/Back/Forward}) \\
    \addlinespace
    \path{artesao-de-skills} & Leitura/Escrita & Demandas pedag{\'o}gicas e regras de schema & M{\'o}dulo \texttt{SKILL.md} v{\'a}lido e auditado \\
    \bottomrule
  \end{tabularx}
\end{table}

\section{O Cat{\'a}logo das 62 Skills do Acervo}

As 62 skills encontram-se agrupadas em cinco categorias operacionais.

\subsection{Categoria 1: N{\'i}veis de Ensino (5 Skills)}

\begin{table}[htbp]
  \centering
  \small
  \caption{Skills de N{\'i}veis de Ensino}
  \label{tab:cat1-niveis}
  \begin{tabularx}{\linewidth}{>{\raggedright\arraybackslash\ttfamily\small}p{5.4cm}X}
    \toprule
    \textbf{\textnormal{\textbf{Skill (Slug)}}} & \textbf{Objetivo e Escopo Pedag{\'o}gico} \\
    \midrule
    \path{ia-educacao-infantil} & Orienta o uso conforme o Referencial MEC: veda telas e intera{\c c}{\~a}o direta de crian{\c c}as; foca no apoio {\`a} gest{\~a}o docente e planejamento l{\'u}dico. \\
    \addlinespace
    \path{ia-educacao-basica} & Integra{\c c}{\~a}o gradual nos anos fundamentais com salvaguardas de privacidade e letramento inicial. \\
    \addlinespace
    \path{ia-educacao-ensino-medio} & Letramento em IA, impactos no mundo do trabalho, c{\'o}digo introdut{\'o}rio e pensamento cr{\'i}tico. \\
    \addlinespace
    \path{ia-educacao-profissional-tecnologica} & Integra{\c c}{\~a}o na EPT, alinhando compet{\^e}ncias de IA {\`a}s demandas dos setores produtivos e ind{\'u}stria 4.0. \\
    \addlinespace
    \path{ia-educacao-superior} & Incorpora{\c c}{\~a}o na gradua{\c c}{\~a}o e p{\'o}s-gradua{\c c}{\~a}o, abrangendo ensino, pesquisa e extens{\~a}o universit{\'a}ria. \\
    \bottomrule
  \end{tabularx}
\end{table}

\subsection{Categoria 2: Forma{\c c}{\~a}o Docente (18 Skills)}

\begin{table}[htbp]
  \centering
  \small
  \caption{Skills de Forma{\c c}{\~a}o Docente}
  \label{tab:cat2-formacao}
  \begin{tabularx}{\linewidth}{>{\raggedright\arraybackslash\ttfamily\small}p{5.4cm}X}
    \toprule
    \textbf{\textnormal{\textbf{Skill (Slug)}}} & \textbf{Objetivo e Escopo Pedag{\'o}gico} \\
    \midrule
    \path{ia-educacao-fundamentos} & Princ{\'i}pios essenciais de IA generativa para educadores alinhados ao Referencial MEC. \\
    \addlinespace
    \path{ia-educacao-formacao-inicial-docente} & Orienta{\c c}{\~a}o para licenciaturas, capacitando futuros professores a mediar a IA. \\
    \addlinespace
    \path{ia-educacao-planejamento-didatico} & Planejamento de aulas completas, sequ{\^e}ncias did{\'a}ticas e engajamento ativo. \\
    \addlinespace
    \path{ia-educacao-planejamento-reverso} & Framework \emph{Backward Design} (Wiggins \& McTighe): resultados, evid{\^e}ncias e plano. \\
    \addlinespace
    \path{ia-educacao-aprendizagem-ativa} & Integra{\c c}{\~a}o de IA com m{\'e}todos ativos, superando o consumo passivo. \\
    \addlinespace
    \path{ia-educacao-bloom} & Alinhamento {\`a} Taxonomia Revisada de Bloom com {\^e}nfase em processos superiores. \\
    \addlinespace
    \path{ia-educacao-metacognicao} & Desenvolvimento da auto-observa{\c c}{\~a}o cognitiva e autorregula{\c c}{\~a}o com IA. \\
    \addlinespace
    \path{ia-educacao-pensamento-critico} & Desenho de tarefas onde a IA atua como est{\'i}mulo e contraponto argumentativo. \\
    \addlinespace
    \path{ia-educacao-pbl} & Aprendizagem Baseada em Problemas e Projetos com co-design inteligente. \\
    \addlinespace
    \path{ia-educacao-sala-invertida} & Estrutura{\c c}{\~a}o dos momentos pr{\'e}-classe e aplica{\c c}{\~a}o ativa presencial. \\
    \addlinespace
    \path{ia-educacao-peer-instruction} & Metodologia de Mazur: formula{\c c}{\~a}o de \emph{ConcepTests} e modera{\c c}{\~a}o por pares. \\
    \addlinespace
    \path{ia-educacao-tbl} & \emph{Team-Based Learning}: elabora{\c c}{\~a}o de iRAT, tRAT e resolu{\c c}{\~a}o em equipe. \\
    \addlinespace
    \path{ia-educacao-estudo-de-caso} & Reda{\c c}{\~a}o de estudos de caso aut{\^e}nticos, complexos e multidisciplinares. \\
    \addlinespace
    \path{ia-educacao-simulacao} & Simula{\c c}{\~o}es de cen{\'a}rios reais e din{\^a}micas de \emph{role-playing} com IA. \\
    \addlinespace
    \path{ia-educacao-debate} & Debates acad{\^e}micos estruturados: identifica{\c c}{\~a}o de fal{\'a}cias e contra-argumenta{\c c}{\~a}o. \\
    \addlinespace
    \path{ia-educacao-facilitacao} & T{\'e}cnicas de media{\c c}{\~a}o de oficinas e trabalhos colaborativos mediados por IA. \\
    \addlinespace
    \path{ia-educacao-aprendizagem-servico} & Integra{\c c}{\~a}o de extens{\~a}o universit{\'a}ria e demandas sociais com solu{\c c}{\~o}es de IA. \\
    \addlinespace
    \path{ia-educacao-interdisciplinaridade} & Articula{\c c}{\~a}o de m{\'u}ltiplas {\'a}reas do saber em projetos curriculares integrados. \\
    \bottomrule
  \end{tabularx}
\end{table}

\subsection{Categoria 3: {\'E}tica e Governan{\c c}a (8 Skills)}

\begin{table}[htbp]
  \centering
  \small
  \caption{Skills de {\'E}tica e Governan{\c c}a}
  \label{tab:cat3-etica}
  \begin{tabularx}{\linewidth}{>{\raggedright\arraybackslash\ttfamily\small}p{5.4cm}X}
    \toprule
    \textbf{\textnormal{\textbf{Skill (Slug)}}} & \textbf{Objetivo e Escopo Pedag{\'o}gico} \\
    \midrule
    \path{ia-educacao-etica} & Fundamenta{\c c}{\~a}o {\'e}tica geral, autonomia humana e justi{\c c}a distributiva em IA. \\
    \addlinespace
    \path{ia-educacao-governanca-dados} & Estrutura{\c c}{\~a}o de pol{\'i}ticas de governan{\c c}a de dados e consentimento discente. \\
    \addlinespace
    \path{ia-educacao-supervisao-humana} & Mecanismos \emph{human-in-the-loop}, vedando avalia{\c c}{\~o}es puramente algor{\'i}tmicas. \\
    \addlinespace
    \path{ia-educacao-transparencia-explicabilidade} & Requisitos de explicabilidade de modelos adotados no ambiente universit{\'a}rio. \\
    \addlinespace
    \path{ia-educacao-vieses} & Identifica{\c c}{\~a}o e mitiga{\c c}{\~a}o de vieses culturais e preconceitos algor{\'i}tmicos. \\
    \addlinespace
    \path{ia-educacao-impacto-algoritmico} & Condu{\c c}{\~a}o de Avalia{\c c}{\~o}es de Impacto Algor{\'i}tmico (AIA) institucionais. \\
    \addlinespace
    \path{ia-educacao-seguranca-digital} & Higiene cibern{\'e}tica e prote{\c c}{\~a}o de dados escolares contra incidentes de seguran{\c c}a. \\
    \addlinespace
    \path{ia-educacao-integridade-academica} & Diretrizes de combate ao pl{\'a}gio e contrato pedag{\'o}gico de autoria respons{\'a}vel. \\
    \bottomrule
  \end{tabularx}
\end{table}

\subsection{Categoria 4: Inclus{\~a}o e Equidade (5 Skills)}

\begin{table}[htbp]
  \centering
  \small
  \caption{Skills de Inclus{\~a}o e Equidade}
  \label{tab:cat4-inclusao}
  \begin{tabularx}{\linewidth}{>{\raggedright\arraybackslash\ttfamily\small}p{5.4cm}X}
    \toprule
    \textbf{\textnormal{\textbf{Skill (Slug)}}} & \textbf{Objetivo e Escopo Pedag{\'o}gico} \\
    \midrule
    \path{ia-educacao-acessibilidade-inclusao} & Recursos de acessibilidade, leitores de tela e tecnologias assistivas com IA. \\
    \addlinespace
    \path{ia-educacao-dua} & Opera{\c c}{\~a}o do Desenho Universal para a Aprendizagem (engajamento, representa{\c c}{\~a}o, a{\c c}{\~a}o). \\
    \addlinespace
    \path{ia-educacao-equidade-digital} & Enfrentamento de desigualdades materiais de conectividade e acesso a computadores. \\
    \addlinespace
    \path{ia-educacao-letramento-dados} & Alfabetiza{\c c}{\~a}o cr{\'i}tica em estat{\'i}stica e visualiza{\c c}{\~a}o de dados educacionais. \\
    \addlinespace
    \path{ia-educacao-permanencia} & Uso preditivo e preventivo de dados acad{\^e}micos para mitigar a evas{\~a}o discente. \\
    \bottomrule
  \end{tabularx}
\end{table}

\subsection{Categoria 5: Ferramentas Pr{\'a}ticas (26 Skills)}

\begin{table}[htbp]
  \centering
  \small
  \caption{Skills de Ferramentas Pr{\'a}ticas (Parte 1)}
  \label{tab:cat5-ferramentas-1}
  \begin{tabularx}{\linewidth}{>{\raggedright\arraybackslash\ttfamily\small}p{5.4cm}X}
    \toprule
    \textbf{\textnormal{\textbf{Skill (Slug)}}} & \textbf{Objetivo e Escopo Pedag{\'o}gico} \\
    \midrule
    \path{aias-consultant} & Consultoria especializada nos 5 n{\'i}veis da Escala AIAS e pol{\'i}ticas de avalia{\c c}{\~a}o. \\
    \addlinespace
    \path{ia-educacao-avaliacao} & Redesenho de instrumentos avaliativos para cen{\'a}rios com presen{\c c}a de IA. \\
    \addlinespace
    \path{ia-educacao-avaliacao-competencia} & Avalia{\c c}{\~a}o baseada em compet{\^e}ncias: rubricas de profici{\^e}ncia e evid{\^e}ncias. \\
    \addlinespace
    \path{ia-educacao-avaliacao-diagnostica} & Mapeamento de pr{\'e}-requisitos conceituais e lacunas no in{\'i}cio da disciplina. \\
    \addlinespace
    \path{ia-educacao-avaliacao-grupo} & Avalia{\c c}{\~a}o do trabalho colaborativo com foco na contribui{\c c}{\~a}o individual. \\
    \addlinespace
    \path{ia-educacao-avaliacao-oral} & Roteiros, rubricas e protocolos para semin{\'a}rios e defesas de projetos. \\
    \addlinespace
    \path{ia-educacao-avaliacao-projeto} & Avalia{\c c}{\~a}o de projetos semestrais e TCCs com entregas intermedi{\'a}rias. \\
    \addlinespace
    \path{ia-educacao-autoavaliacao} & Instrumentos reflexivos para autoavalia{\c c}{\~a}o e autorregula{\c c}{\~a}o discente. \\
    \addlinespace
    \path{ia-educacao-banco-questoes} & Gest{\~a}o de bancos de quest{\~o}es categorizadas por n{\'i}vel de Bloom e blueprints. \\
    \addlinespace
    \path{ia-educacao-mcq} & Design e valida{\c c}{\~a}o de itens de m{\'u}ltipla escolha com distratores calibrados. \\
    \addlinespace
    \path{ia-educacao-portfolio} & Design de portf{\'o}lios avaliativos e di{\'a}rios de bordo com reflex{\~a}o metacognitiva. \\
    \addlinespace
    \path{ia-educacao-rubrica} & Constru{\c c}{\~a}o de rubricas anal{\'i}ticas com r{\'o}tulos neutros e coer{\^e}ncia AIAS. \\
    \addlinespace
    \path{ia-educacao-feedback} & Design de feedback formativo no modelo \emph{Feed Up/Back/Forward}. \\
    \bottomrule
  \end{tabularx}
\end{table}

\begin{table}[htbp]
  \centering
  \small
  \caption{Skills de Ferramentas Pr{\'a}ticas (Parte 2)}
  \label{tab:cat5-ferramentas-2}
  \begin{tabularx}{\linewidth}{>{\raggedright\arraybackslash\ttfamily\small}p{5.4cm}X}
    \toprule
    \textbf{\textnormal{\textbf{Skill (Slug)}}} & \textbf{Objetivo e Escopo Pedag{\'o}gico} \\
    \midrule
    \path{ia-educacao-rascunho} & T{\'e}cnica \emph{Chain of Draft} (CoD): gera{\c c}{\~a}o de esbo{\c c}os intermedi{\'a}rios para reflex{\~a}o. \\
    \addlinespace
    \path{ia-educacao-verificacao} & T{\'e}cnica \emph{Chain of Verification} (CoVe): auditoria independente contra alucina{\c c}{\~o}es. \\
    \addlinespace
    \path{ia-educacao-personalizacao} & Trilhas de aprendizagem diferenciadas conforme o ritmo e dificuldades do aluno. \\
    \addlinespace
    \path{ia-educacao-sti} & Sistemas Tutoriais Inteligentes (STIs) para resolu{\c c}{\~a}o assistida de problemas. \\
    \addlinespace
    \path{ia-educacao-ia-desplugada} & Atividades sem dispositivos digitais para ensino de l{\'o}gica de IA e {\'e}tica. \\
    \addlinespace
    \path{ia-educacao-sandbox} & Ambientes controlados para testagem segura de novas aplica{\c c}{\~o}es educacionais. \\
    \addlinespace
    \path{ia-educacao-gestao} & Compet{\^e}ncias de gest{\~a}o acad{\^e}mica para coordenadores e dirigentes de curso. \\
    \addlinespace
    \path{ia-educacao-ecossistema-inovacao} & Parcerias universidade-empresa e laborat{\'o}rios de inova{\c c}{\~a}o pedag{\'o}gica. \\
    \addlinespace
    \path{ia-educacao-contratacao} & Crit{\'e}rios t{\'e}cnicos, jur{\'i}dicos e pedag{\'o}gicos para licita{\c c}{\~a}o de solu{\c c}{\~o}es de IA. \\
    \addlinespace
    \path{ia-educacao-design-problema} & Formula{\c c}{\~a}o de desafios aut{\^e}nticos e situa{\c c}{\~o}es-gatilho para disciplinas STHEM. \\
    \addlinespace
    \path{ia-educacao-pesquisa} & Metodologia de pesquisa cient{\'i}fica, revis{\~a}o integrativa e an{\'a}lise com IA. \\
    \addlinespace
    \path{ia-educacao-escrita} & Escrita acad{\^e}mica assistida por IA garantindo preserva{\c c}{\~a}o da voz autoral. \\
    \addlinespace
    \path{ia-educacao-visualizacao-dados} & Comunica{\c c}{\~a}o cient{\'i}fica visual, gr{\'a}ficos informativos e \emph{storytelling} com dados. \\
    \bottomrule
  \end{tabularx}
\end{table}

\chapter{Casos de Uso Iniciais}
\label{cap:casos-de-uso}

\begin{objetivos}
\begin{itemize}
  \item Reproduzir, no seu contexto, os quatro roteiros pr{\'a}ticos apresentados (planejamento, avalia{\c c}{\~a}o, feedback e sala de aula ativa);
  \item Adaptar os exemplos de \emph{prompts} {\`a} sua {\'a}rea do conhecimento e {\`a} sua disciplina;
  \item Aplicar o checklist de valida{\c c}{\~a}o docente antes de disponibilizar materiais gerados com IA.
\end{itemize}
\end{objetivos}

\section{Introdu{\c c}{\~a}o aos Cen{\'a}rios Pr{\'a}ticos}

Este cap{\'i}tulo foi desenhado para servir como um guia de aplica{\c c}{\~a}o r{\'a}pida. Nele, apresentamos quatro casos de uso t{\'i}picos enfrentados por professores do ensino superior e da educa{\c c}{\~a}o profissional, detalhando as entradas do docente, as chamadas de agentes e os artefatos finais resultantes. A Tabela~\ref{tab:casos-indice} permite localizar rapidamente o caso mais pr{\'o}ximo da sua necessidade.

\begin{table}[htbp]
  \centering
  \small
  \caption{Como Escolher o Caso de Uso Adequado}
  \label{tab:casos-indice}
  \begin{tabularx}{\linewidth}{p{3.2cm}X}
    \toprule
    \textbf{Necessidade Docente} & \textbf{Caso e {\'A}rea Ilustrativa} \\
    \midrule
    Estruturar uma disciplina inteira, definir n{\'i}veis AIAS e cronograma semestral & Caso~1 -- \emph{Estruturas de Dados e Algoritmos} (Ci{\^e}ncia da Computa{\c c}{\~a}o / Engenharias) \\
    \addlinespace
    Elaborar uma avalia{\c c}{\~a}o formativa completa (enunciado, rubrica e quest{\~o}es objetivas) & Caso~2 -- \emph{Banco de Dados} (Ci{\^e}ncia da Computa{\c c}{\~a}o / Engenharias) \\
    \addlinespace
    Gerar feedback formativo inclusivo para estudantes com dislexia ou TDAH & Caso~3 -- Qualquer {\'a}rea com entrega escrita ou c{\'o}digo-fonte \\
    \addlinespace
    Transformar uma aula expositiva em atividade ativa de \emph{Peer Instruction} & Caso~4 -- \emph{Termodin{\^a}mica} (F{\'i}sica / Engenharias) \\
    \bottomrule
  \end{tabularx}
\end{table}

Cada caso segue uma estrutura fixa -- \emph{Contexto e Desafio}, \emph{Execu{\c c}{\~a}o no OpenCode} e \emph{Artefato Entregue} -- facilitando a leitura transversal e a adapta{\c c}{\~a}o para outras disciplinas.

\begin{atencaopedagogica}
\textbf{Como invocar skills e subagentes}: O OpenCode adota duas sintaxes distintas de acionamento. \textbf{Skills} s{\~a}o invocadas com a barra no campo de mensagem, na forma \texttt{/ia-educacao-rubrica} (ou \texttt{/aias-consultant}); alternativamente, podem ser mencionadas no texto do \emph{prompt} pelo slug. \textbf{Subagentes} s{\~a}o acionados com o arroba, na forma \texttt{@planejador-pedagogico} (ou \texttt{@coach-de-feedback-formativo}), ou pela sele{\c c}{\~a}o no menu suspenso da interface. Nos exemplos deste cap{\'i}tulo, os \emph{prompts} seguem essas conven{\c c}{\~o}es: \texttt{@} para o subagente e slug simples (ou \texttt{/} quando isolado) para a skill. Os conceitos de skill e de subagente, bem como os cat{\'a}logos completos, encontram-se respectivamente nos Cap{\'i}tulos~\ref{cap:skills-opencode} e~\ref{cap:agentes-skills}.
\end{atencaopedagogica}

\section{Caso 1: Planejando uma Disciplina do Zero com \emph{Backward Design}}

\subsection{Contexto e Desafio}
O professor assume a disciplina de \emph{Estruturas de Dados e Algoritmos} (60 horas) e deseja estruturar um plano de ensino inovador, alinhado ao \emph{Referencial do MEC} \cite{brasil2026referencial}, com metodologias ativas e enquadramento intencional na Escala AIAS.

\subsection{Execu{\c c}{\~a}o no OpenCode}
No terminal ou na interface gr{\'a}fica, o professor invoca o subagente consultivo:

\begin{lstlisting}[language=bash]
opencode
\end{lstlisting}

\begin{tcolorbox}[colback=white, colframe=boxborder, arc=1mm, title={Comando Enviado pelo Professor}]
\small
``Atue como o subagente \path{@planejador-pedagogico}. Analise a ementa de \emph{Estruturas de Dados} (listas, pilhas, filas, {\'a}rvores bin{\'a}rias, grafos e algoritmos de ordena{\c c}{\~a}o). Proponha um planejamento reverso de 15 semanas com tr{\^e}s momentos avaliativos formativos e especifique o n{\'i}vel AIAS recomendado para cada unidade tem{\'a}tica.''
\end{tcolorbox}

No OpenCode, a forma can{\^o}nica de acionar o subagente {\'e} a men{\c c}{\~a}o pelo nome no pr{\'o}prio texto do \emph{prompt}, na sintaxe ``\path{@planejador-pedagogico}'', ou a sele{\c c}{\~a}o do agente no menu suspenso da interface gr{\'a}fica.

\subsection{Artefato Entregue}
O subagente aplica os frameworks (AIAS, Bloom e \emph{Backward Design}) e indica as skills do reposit{\'o}rio que apoiam a implementa{\c c}{\~a}o do plano (incluindo \path{aias-consultant}, \path{ia-educacao-avaliacao}, \path{ia-educacao-rubrica} e \path{ia-educacao-planejamento-didatico}), gerando:
\begin{itemize}
  \item \textbf{Resultados Desejados}: Compet{\^e}ncias centrais focadas nos n{\'i}veis \emph{Analisar} e \emph{Criar} (avaliar complexidade assint{\'o}tica e projetar estruturas adequadas para problemas reais);
  \item \textbf{Enquadramento AIAS Progressivo}:
    \begin{itemize}
      \item \emph{Semanas 1 a 4 (Ponteiros e Listas)}: N{\'i}vel AIAS~1 (Sem IA), garantindo fixa{\c c}{\~a}o manual de refer{\^e}ncias de mem{\'o}ria;
      \item \emph{Semanas 5 a 10 ({\'A}rvores e Algoritmos)}: N{\'i}vel AIAS~2 (Planejamento Assistido) e AIAS~3 (Colabora{\c c}{\~a}o na depura{\c c}{\~a}o de casos de teste);
      \item \emph{Semanas 11 a 15 (Projeto com Grafos)}: N{\'i}vel AIAS~4 (IA Integral), onde a equipe projeta um sistema de rotas usando IA para co-gerar implementa{\c c}{\~o}es sob rigorosa auditoria discente.
    \end{itemize}
  \item \textbf{Plano de Aulas}: Cronograma de 15 semanas articulado com desafios pr{\'a}ticos em laborat{\'o}rio.
\end{itemize}

\section{Caso 2: Elaborando uma Avalia{\c c}{\~a}o Formativa com Rubrica e MCQs}

\subsection{Contexto e Desafio}
Para o fechamento da segunda unidade de uma disciplina de \emph{Banco de Dados}, o docente precisa de um instrumento avaliativo aut{\^e}ntico com declara{\c c}{\~a}o formal AIAS e uma rubrica anal{\'i}tica detalhada.

\subsection{Execu{\c c}{\~a}o no OpenCode}
\begin{tcolorbox}[colback=white, colframe=boxborder, arc=1mm, title={Comando Enviado pelo Professor}]
\small
``Atue como o subagente \path{@construtor-de-avaliacoes}. Crie uma avalia{\c c}{\~a}o formativa sobre \emph{Modelagem Relacional e Normaliza{\c c}{\~a}o (1FN a 3FN)} para um cen{\'a}rio de com{\'e}rcio eletr{\^o}nico. A atividade ser{\'a} realizada no n{\'i}vel AIAS~2. Em seguida, elabore a rubrica anal{\'i}tica em 4 n{\'i}veis com a skill \path{ia-educacao-rubrica} e adicione duas quest{\~o}es de m{\'u}ltipla escolha diagn{\'o}sticas com distratores baseados em equ{\'i}vocos conceituais t{\'i}picos usando a skill \path{ia-educacao-mcq}.''
\end{tcolorbox}

\subsection{Artefato Entregue}
O subagente gera um arquivo Markdown/LaTeX contendo:
\begin{enumerate}
  \item \textbf{Cabe{\c c}alho Institucional com AIAS}: T{\'i}tulo, instru{\c c}{\~o}es de entrega e a declara{\c c}{\~a}o expl{\'i}cita: ``\emph{Esta atividade adota o N{\'i}vel AIAS 2 -- Planejamento Assistido por IA. Voc{\^e} pode debater ideias de tabelas com a ferramenta, mas os diagramas l{\'o}gicos e as senten{\c c}as DDL devem ser integralmente redigidos por voc{\^e}}'';
  \item \textbf{Desafio Aut{\^e}ntico}: Um cen{\'a}rio real de um e-commerce com anomalias de inser{\c c}{\~a}o e atualiza{\c c}{\~a}o propositais para o discente normalizar;
  \item \textbf{Rubrica Anal{\'i}tica em 4 N{\'i}veis}:
    \begin{itemize}
      \item Crit{\'e}rio 1: Identifica{\c c}{\~a}o de Chaves e Depend{\^e}ncias Funcionais;
      \item Crit{\'e}rio 2: Aplica{\c c}{\~a}o das Formas Normais (1FN a 3FN);
      \item Crit{\'e}rio 3: Integridade Referencial e Restri{\c c}{\~o}es;
      \item Crit{\'e}rio 4: Transpar{\^e}ncia no Uso de IA (coer{\^e}ncia com AIAS~2).
    \end{itemize}
  \item \textbf{Dois Itens MCQs Diagn{\'o}sticos}: Cada distrator com uma nota explicativa indicando exatamente a falha de racioc{\'i}nio do aluno que escolheu aquela alternativa (e.g., confundir depend{\^e}ncia parcial com transitiva).
\end{enumerate}

\section{Caso 3: Ciclo de Feedback Formativo e Inclus{\~a}o}

\subsection{Contexto e Desafio}
O professor recebeu a entrega de um estudante de gradua{\c c}{\~a}o com diagn{\'o}stico de dislexia. O trabalho apresenta excelentes solu{\c c}{\~o}es algor{\'i}tmicas, mas o relat{\'o}rio textual possui estrutura{\c c}{\~a}o truncada e falta a an{\'a}lise de casos de teste de borda.

\subsection{Execu{\c c}{\~a}o no OpenCode}
\begin{tcolorbox}[colback=white, colframe=boxborder, arc=1mm, title={Comando Enviado pelo Professor}]
\small
``Atue como o subagente \path{@coach-de-feedback-formativo}. Analise a submiss{\~a}o anexa em rela{\c c}{\~a}o {\`a} nossa rubrica. O estudante possui laudo de dislexia. Elabore um feedback formativo nas tr{\^e}s dimens{\~o}es (\emph{Feed Up}, \emph{Feed Back}, \emph{Feed Forward}) aplicando as diretrizes de DUA da skill \path{ia-educacao-dua}: use frases curtas em ordem direta, listas pontuadas, limite as a{\c c}{\~o}es de melhoria a duas prioridades e adicione uma pergunta ativadora da metacogni{\c c}{\~a}o.''
\end{tcolorbox}

\subsection{Artefato Entregue}
O subagente gera um parecer acolhedor, objetivo e isento de sobrecarga textual:
\begin{itemize}
  \item \textbf{Feed Up}: Relembra a meta principal (implementar o algoritmo de Dijkstra com fila de prioridades em complexidade $O(E \log V)$);
  \item \textbf{Feed Back}: Destaca em primeiro lugar o acerto do c{\'o}digo (estrutura de dados correta e boa modulariza{\c c}{\~a}o); aponta em seguida a aus{\^e}ncia de testes com grafos desconexos;
  \item \textbf{Feed Forward}: Duas micro-a{\c c}{\~o}es imediatas (passo 1: adicionar uma condi{\c c}{\~a}o de parada para v{\'e}rtices inalcan{\c c}{\~a}veis; passo 2: testar uma entrada com 3 v{\'e}rtices isolados);
  \item \textbf{Pergunta Metacognitiva}: ``Como voc{\^e} explicaria a l{\'o}gica de parada do algoritmo para um colega sem usar jarg{\~o}es t{\'e}cnicos?''
\end{itemize}

\section{Caso 4: Din{\^a}mica Ativa de Sala Invertida e \emph{Peer Instruction}}

\subsection{Contexto e Desafio}
O docente quer transformar uma aula te{\'o}rica expositiva sobre \emph{Termodin{\^a}mica (Segunda Lei)} em uma sess{\~a}o ativa presencial de 100 minutos baseada em discuss{\~a}o por pares.

\subsection{Execu{\c c}{\~a}o no OpenCode}
\begin{tcolorbox}[colback=white, colframe=boxborder, arc=1mm, title={Comando Enviado pelo Professor}]
\small
``Utilize as skills \path{ia-educacao-sala-invertida} e \path{ia-educacao-peer-instruction} para estruturar a aula de Segunda Lei da Termodin{\^a}mica. Preciso do roteiro de leitura pr{\'e}-aula com tr{\^e}s perguntas disparadoras de d{\'u}vida e tr{\^e}s \emph{ConcepTests} conceituais para vota{\c c}{\~a}o presencial com flashcards/celular.''
\end{tcolorbox}

\subsection{Artefato Entregue}
O OpenCode estrutura:
\begin{itemize}
  \item \textbf{Guia Pr{\'e}-Aula (15 minutos discentes)}: Leitura dirigida de duas p{\'a}ginas com perguntas de autorreflex{\~a}o para o aluno trazer anotadas;
  \item \textbf{Din{\^a}mica Presencial}: Divis{\~a}o dos 100 minutos em mini-blocos explicativos de 10 minutos intercalados por \emph{ConcepTests};
  \item \textbf{Quest{\~o}es Conceituais}: Itens qualitativos sem f{\'o}rmulas num{\'e}ricas que for{\c c}am os estudantes a discutir em duplas sobre entropia e irreversibilidade antes da segunda vota{\c c}{\~a}o.
\end{itemize}

\section{Cen{\'a}rios por {\'A}rea STHEM}

As quatro disciplinas ilustrativas dos casos anteriores concentram-se nas {\'a}reas de tecnologia e ci{\^e}ncias exatas. Para evidenciar que o acervo atende a todo o espectro STHEM (Science, Technology, Humanities, Engineering, Mathematics), esta se{\c c}{\~a}o apresenta um cen{\'a}rio adicional por {\'a}rea, com a skill que o OpenCode seleciona espontaneamente ao receber o \emph{prompt}. Os cinco cen{\'a}rios foram validados no OpenCode real: em todos eles, o agente identificou corretamente as skills do acervo a partir da descri{\c c}{\~a}o da tarefa, sem necessidade de citar o nome da skill no \emph{prompt} -- demonstrando que a detec{\c c}{\~a}o sem{\^a}ntica funciona a partir do objetivo pedag{\'o}gico.

\begin{dicaopencode}
\textbf{Valida{\c c}{\~a}o emp{\'i}rica dos roteiros deste cap{\'i}tulo}: os quatro casos de uso (Casos~1 a~4) e os cinco cen{\'a}rios STHEM desta se{\c c}{\~a}o foram executados no OpenCode real (vers{\~a}o 1.18.29, modo headless \texttt{opencode run}). Em todos, as skills citadas foram efetivamente carregadas pela ferramenta \texttt{skill()} e os artefatos esperados foram produzidos (planos, rubricas com r{\'o}tulos can{\^o}nicos, feedbacks tridimensionais e ConcepTests). Um detalhe t{\'e}cnico observado: ao delegar via \texttt{@subagente}, o OpenCode atual exige um identificador de modelo v{\'a}lido no frontmatter do agente (ex.: \texttt{model: opencode/glm-5.3-flash}); o valor \texttt{model: any} do acervo -- um marcador de independ{\^e}ncia de provedor -- n{\~a}o {\'e} interpretado como ID e faz a delega{\c c}{\~a}o direta falhar. Nesse caso, o agente principal assume o papel do subagente e executa as mesmas skills, preservando o resultado; para restaurar a delega{\c c}{\~a}o via Task, basta substituir \texttt{model: any} por um ID de modelo real no arquivo do agente em \texttt{.opencode/agents/}.
\end{dicaopencode}

Os perfis por {\'a}rea seguem a tabela de adapta{\c c}{\~o}es embutida nos subagentes do reposit{\'o}rio, que orientam n{\'i}veis AIAS t{\'i}picos e formatos preferenciais de avalia{\c c}{\~a}o para cada {\'a}rea (Tabela~\ref{tab:sthem-cenarios}).

\begin{table}[htbp]
  \centering
  \small
  \caption{Cen{\'a}rios Validados por {\'A}rea STHEM}
  \label{tab:sthem-cenarios}
  \begin{tabularx}{\linewidth}{p{2.2cm}p{4.6cm}p{3.6cm}X}
    \toprule
    \textbf{{\'A}rea} & \textbf{Cen{\'a}rio Docente} & \textbf{Skills Selecionadas (validado)} & \textbf{N{\'i}vel AIAS T{\'i}pico} \\
    \midrule
    Science & Relat{\'o}rio de laborat{\'o}rio de microbiologia com an{\'a}lise de dados (Biologia) & \path{ia-educacao-escrita}, \path{ia-educacao-avaliacao-projeto} & 1 (prova), 3 (relat{\'o}rio), 4 (dados) \\
    \addlinespace
    Technology & Projeto em equipe de API com IA integral e auditoria discente (Ci. da Computa{\c c}{\~a}o) & \path{ia-educacao-pbl}, \path{ia-educacao-avaliacao-projeto} & 3 (c{\'o}digo assistido), 4 (projeto), 5 (fronteiras) \\
    \addlinespace
    Humanities & Quest{\~o}es discursivas de an{\'a}lise de fontes prim{\'a}rias (Hist{\'o}ria) & \path{ia-educacao-avaliacao} & 2 (pesquisa), 3 (escrita), 5 (co-cria{\c c}{\~a}o) \\
    \addlinespace
    Engineering & Rubrica de memorial de c{\'a}lculo com crit{\'e}rio de transpar{\^e}ncia de IA (Eng. Civil) & \path{ia-educacao-rubrica} & 1 (c{\'a}lculo), 3 (projeto), 4 (otimiza{\c c}{\~a}o) \\
    \addlinespace
    Mathematics & Banco de quest{\~o}es de limites com distratores por equ{\'i}voco (C{\'a}lculo 1) & \path{ia-educacao-banco-questoes}, \path{ia-educacao-mcq} & 1 (dom{\'i}nio), 2 (explora{\c c}{\~a}o), 3 (verifica{\c c}{\~a}o) \\
    \bottomrule
  \end{tabularx}
\end{table}

\subsection{Science: Relat{\'o}rio de Laborat{\'o}rio (Biologia)}

\begin{tcolorbox}[colback=white, colframe=boxborder, arc=1mm, title={Prompt Validado -- Science}]
\small
``Sou professor de Biologia e preciso elaborar o roteiro de um relat{\'o}rio de laborat{\'o}rio de microbiologia que avalie m{\'e}todo cient{\'i}fico e an{\'a}lise de dados, com n{\'i}vel AIAS 3.''
\end{tcolorbox}

\textbf{Resposta do OpenCode (validada)}: o agente seleciona espontaneamente \path{ia-educacao-escrita} (roteiro de relat{\'o}rio t{\'e}cnico, metodologia com IA como apoio sem perder a voz autoral) e \path{ia-educacao-avaliacao-projeto} (rubrica com crit{\'e}rios de an{\'a}lise de dados e marcos de entrega), sinalizando ainda que a declara{\c c}{\~a}o do n{\'i}vel AIAS~3 deve ser calibrada com \path{aias-consultant}.

\subsection{Technology: Projeto Colaborativo com IA Integral (Ci{\^e}ncia da Computa{\c c}{\~a}o)}

\begin{tcolorbox}[colback=white, colframe=boxborder, arc=1mm, title={Prompt Validado -- Technology}]
\small
``Quero planejar uma atividade de projeto onde equipes desenvolvem uma API com assist{\^e}ncia integral de IA (n{\'i}vel AIAS 4), com entregas incrementais e auditoria discente do c{\'o}digo gerado.''
\end{tcolorbox}

\textbf{Resposta do OpenCode (validada)}: o agente seleciona \path{ia-educacao-pbl} (din{\^a}mica de equipes com entregas incrementais) e \path{ia-educacao-avaliacao-projeto} (rubricas e auditoria de projeto com IA, coerentes com o n{\'i}vel AIAS~4).

\subsection{Humanities: An{\'a}lise de Fontes Prim{\'a}rias (Hist{\'o}ria)}

\begin{tcolorbox}[colback=white, colframe=boxborder, arc=1mm, title={Prompt Validado -- Humanities}]
\small
``Elabore tr{\^e}s quest{\~o}es discursivas sobre a Revolu{\c c}{\~a}o Industrial que avaliem an{\'a}lise de fontes prim{\'a}rias.''
\end{tcolorbox}

\textbf{Resposta do OpenCode (validada)}: o agente seleciona \path{ia-educacao-avaliacao} -- e distingue corretamente que, por se tratarem de quest{\~o}es \emph{discursivas} (e n{\~a}o de m{\'u}ltipla escolha), a skill \path{ia-educacao-mcq} n{\~a}o {\'e} a adequada. A arg{\'u}cia de sele{\c c}{\~a}o demonstra que as descri{\c c}{\~o}es das skills guiam a escolha com precis{\~a}o.

\subsection{Engineering: Rubrica de Memorial de C{\'a}lculo (Engenharia Civil)}

\begin{tcolorbox}[colback=white, colframe=boxborder, arc=1mm, title={Prompt Validado -- Engineering}]
\small
``Crie uma rubrica anal{\'i}tica de 4 n{\'i}veis para avaliar o memorial de c{\'a}lculo de estruturas de concreto armado, incluindo crit{\'e}rio de transpar{\^e}ncia no uso de IA (n{\'i}vel AIAS 1 para o c{\'a}lculo).''
\end{tcolorbox}

\textbf{Resposta do OpenCode (validada)}: o agente seleciona \path{ia-educacao-rubrica} -- rubricas anal{\'i}ticas em n{\'i}veis, alinhadas a Bloom, AIAS e DUA, com o crit{\'e}rio de transpar{\^e}ncia no uso de IA. O cen{\'a}rio evidencia a independ{\^e}ncia entre Bloom e AIAS: o c{\'a}lculo pode ser AIAS~1 enquanto a an{\'a}lise cr{\'i}tica do memorial admite n{\'i}veis superiores.

\subsection{Mathematics: Banco de Quest{\~o}es com Distratores (C{\'a}lculo 1)}

\begin{tcolorbox}[colback=white, colframe=boxborder, arc=1mm, title={Prompt Validado -- Mathematics}]
\small
``Monte um banco de 15 quest{\~o}es sobre limites, categorizadas por n{\'i}vel de Bloom, com distratores baseados em equ{\'i}vocos conceituais t{\'i}picos (ex.: confundir limite com valor da fun{\c c}{\~a}o).''
\end{tcolorbox}

\textbf{Resposta do OpenCode (validada)}: o agente seleciona \path{ia-educacao-banco-questoes} (estrutura do banco e blueprint por n{\'i}vel de Bloom) e \path{ia-educacao-mcq} (distratores plaus{\'i}veis ancorados em equ{\'i}vocos reais de racioc{\'i}nio).

\begin{atencaopedagogica}
\textbf{Como reproduzir os cen{\'a}rios}: os cinco \emph{prompts} acima foram testados no OpenCode real e selecionam corretamente as skills sem cit{\'a}-las pelo nome -- basta descrever a tarefa com clareza. O docente pode adaptar o tema (ex.: trocar ``Revolu{\c c}{\~a}o Industrial'' por ``Modernismo brasileiro'') mantendo a estrutura do pedido; a sele{\c c}{\~a}o sem{\^a}ntica da skill permanece a mesma. Para maior determinismo, a skill pode ainda ser citada explicitamente (\texttt{/slug} ou men{\c c}{\~a}o no texto).
\end{atencaopedagogica}

\section{Checklist de Valida{\c c}{\~a}o Docente Pr{\'e}-Aplica{\c c}{\~a}o}

Antes de disponibilizar qualquer material gerado com o aux{\'i}lio do OpenCode aos estudantes, o professor deve realizar a checagem dos cinco pontos da Tabela~\ref{tab:checklist-docente}.

\begin{table}[htbp]
  \centering
  \small
  \caption{Checklist de Valida{\c c}{\~a}o Docente}
  \label{tab:checklist-docente}
  \begin{tabularx}{\linewidth}{clp{3.5cm}X}
    \toprule
    \textbf{Item} & \textbf{Crit{\'e}rio de Controle} & \textbf{Pergunta de Checagem} & \textbf{A{\c c}{\~a}o em Caso de N{\~a}o} \\
    \midrule
    \textbf{1} & Coer{\^e}ncia Pedag{\'o}gica & A tarefa espelha os objetivos reais do curso? & Ajuste os verbos com a skill \path{ia-educacao-bloom}. \\
    \addlinespace
    \textbf{2} & Declara{\c c}{\~a}o AIAS & O n{\'i}vel de IA est{\'a} expl{\'i}cito no enunciado? & Adicione a cl{\'a}usula de integridade com \path{aias-consultant}. \\
    \addlinespace
    \textbf{3} & Auditoria Factual & As f{\'o}rmulas e c{\'o}digos foram rodados? & Execute o teste real e audite com \path{ia-educacao-verificacao}. \\
    \addlinespace
    \textbf{4} & Acessibilidade DUA & A linguagem {\'e} clara e inclusiva? & Simplifique par{\'a}grafos densos com \path{ia-educacao-dua}. \\
    \addlinespace
    \textbf{5} & Autonomia Humana & A avalia{\c c}{\~a}o final depende do professor? & Garanta que o parecer passe pelo seu crivo cr{\'i}tico. \\
    \bottomrule
  \end{tabularx}
\end{table}


% --- Bibliografia ABNT ---
\backmatter
\bibliography{referencias}

\end{document}
